% =============================================================================
% Release 2 — Industrialisation et déploiement de la plateforme TicketManager
% =============================================================================
%
% Inclusion dans le rapport principal (mémoire / chapitre) :
%
%   \chapter{Release 2 : Industrialisation et déploiement de la plateforme}
%   \graphicspath{{docs/release-2/img/}{docs/release-1/img/}}
%   % =============================================================================
% Release 2 — Industrialisation et déploiement de la plateforme TicketManager
% =============================================================================
%
% Inclusion dans le rapport principal (mémoire / chapitre) :
%
%   \chapter{Release 2 : Industrialisation et déploiement de la plateforme}
%   \graphicspath{{docs/release-2/img/}{docs/release-1/img/}}
%   % =============================================================================
% Release 2 — Industrialisation et déploiement de la plateforme TicketManager
% =============================================================================
%
% Inclusion dans le rapport principal (mémoire / chapitre) :
%
%   \chapter{Release 2 : Industrialisation et déploiement de la plateforme}
%   \graphicspath{{docs/release-2/img/}{docs/release-1/img/}}
%   % =============================================================================
% Release 2 — Industrialisation et déploiement de la plateforme TicketManager
% =============================================================================
%
% Inclusion dans le rapport principal (mémoire / chapitre) :
%
%   \chapter{Release 2 : Industrialisation et déploiement de la plateforme}
%   \graphicspath{{docs/release-2/img/}{docs/release-1/img/}}
%   \input{docs/release-2/release-2}
%
% Convention : les \caption des tableaux sont placés AU-DESSUS du tabular.
% Les figures de réalisation sont citées dans le texte via \ref{fig:...}.
%
% =============================================================================

\section*{Release 2 : Industrialisation et déploiement de la plateforme}
\addcontentsline{toc}{section}{Release 2 : Industrialisation et déploiement de la plateforme}

La \textbf{Release~2} prolonge le socle applicatif livré en Release~1 (conception et
développement de \textit{TicketManager}). Elle vise à industrialiser le cycle de vie de la
plateforme sur \textbf{Azure Stack Hub} : provisionnement reproductible de l'infrastructure,
chaîne d'intégration et de déploiement continus (CI/CD), déploiement conteneurisé sur
Kubernetes et supervision de l'environnement de production.

\paragraph{Organisation en trois sprints.}
Chaque sprint suit la même trame pédagogique que la Release~1 :
\begin{enumerate}
  \item \textbf{objectifs} fixés et périmètre ;
  \item \textbf{backlog} planifié (user stories) ;
  \item \textbf{analyse} (conception technique, architecture) ;
  \item \textbf{réalisation} (mise en œuvre, captures d'écran, validation).
\end{enumerate}

\paragraph{Codification des priorités (backlog).}
Dans les tableaux de backlog, la colonne \textbf{Priorité} utilise les codes suivants :
\textbf{H} (haute), \textbf{M} (moyenne) et \textbf{B} (basse).

Le tableau~\ref{tab:r2-overview} synthétise les trois sprints de la release.

\begin{table}[H]
\centering
\caption{Vue d'ensemble de la Release~2}
\label{tab:r2-overview}
\footnotesize
\renewcommand{\arraystretch}{1.15}
\setlength{\tabcolsep}{6pt}
\begin{tabular}{|p{1.2cm}|p{4.2cm}|p{7.8cm}|}
\hline
\textbf{Sprint} & \textbf{Intitulé} & \textbf{Livrables principaux} \\
\hline
Sprint~1 & Mise en place de l'infrastructure &
  Terraform (VNet, NSG, VMs), AKS Engine (cluster K8s), plan d'adressage \\
\hline
Sprint~2 & Intégration continue et déploiement continu (CI/CD) &
  Pipeline Azure DevOps, SonarQube, Docker, Nexus, Trivy, Terraform automatisé \\
\hline
Sprint~3 & Monitoring et supervision de la plateforme &
  Déploiement K8s (CD), MongoDB in-cluster, Prometheus, Grafana (Ansible) \\
\hline
\end{tabular}
\end{table}

\paragraph{Périmètre par rapport à la Release~1.}
La Release~1 couvre le développement fonctionnel (NestJS, Next.js, MongoDB locale,
Cloudinary, notifications WebSocket). La Release~2 ne modifie pas le métier applicatif :
elle prépare et automatise l'environnement cloud, la livraison des images Docker et le
déploiement des composants \texttt{ticket-backend}, \texttt{ticket-frontend} et
\texttt{mongodb} dans le namespace Kubernetes \texttt{ticketmanager}.

\paragraph{Fichiers sources du dépôt.}
\begin{itemize}
  \item Infrastructure : répertoire \texttt{terraform/} ;
  \item Pipeline : \texttt{azure-pipelines.yml} ;
  \item Déploiement Kubernetes : répertoire \texttt{k8s/} ;
  \item Automatisation : répertoire \texttt{ansible/} ;
  \item Documentation Release~1 (socle applicatif) : \texttt{docs/release-1/}.
\end{itemize}

% --- Sprint 1 : Infrastructure ------------------------------------------------
\input{docs/release-2/sprint-1-infrastructure}

% --- Sprint 2 : CI/CD ---------------------------------------------------------
\input{docs/release-2/sprint-2-cicd}

% --- Sprint 3 : Monitoring et CD ----------------------------------------------
\input{docs/release-2/sprint-3-monitoring-cd}

% --- Bilan Release 2 ----------------------------------------------------------
\subsection{Bilan de la Release~2}

À l'issue des trois sprints, la plateforme \textit{TicketManager} dispose d'une
infrastructure documentée et reproductible sur Azure Stack Hub, d'une chaîne CI/CD
automatisée (compilation, qualité, conteneurisation, publication Nexus, provisionnement
IaC) et d'un déploiement applicatif supervisé (pods Kubernetes, Prometheus, Grafana).

Les indicateurs de réussite retenus pour cette release sont :
\begin{itemize}
  \item exécution réussie du pipeline Azure DevOps sur la branche \texttt{main} ;
  \item pods du namespace \texttt{ticketmanager} en état \texttt{Running} ;
  \item accessibilité de l'interface web et de l'API depuis le réseau autorisé ;
  \item cibles Prometheus en état \texttt{UP} et dashboards Grafana opérationnels.
\end{itemize}

Les évolutions ultérieures (hors Release~2) peuvent porter sur le renforcement des tests
automatisés dans le pipeline, la réactivation des stages Snyk/Trivy, l'exposition via
Ingress Controller et la haute disponibilité multi-réplicas.

%
% Convention : les \caption des tableaux sont placés AU-DESSUS du tabular.
% Les figures de réalisation sont citées dans le texte via \ref{fig:...}.
%
% =============================================================================

\section*{Release 2 : Industrialisation et déploiement de la plateforme}
\addcontentsline{toc}{section}{Release 2 : Industrialisation et déploiement de la plateforme}

La \textbf{Release~2} prolonge le socle applicatif livré en Release~1 (conception et
développement de \textit{TicketManager}). Elle vise à industrialiser le cycle de vie de la
plateforme sur \textbf{Azure Stack Hub} : provisionnement reproductible de l'infrastructure,
chaîne d'intégration et de déploiement continus (CI/CD), déploiement conteneurisé sur
Kubernetes et supervision de l'environnement de production.

\paragraph{Organisation en trois sprints.}
Chaque sprint suit la même trame pédagogique que la Release~1 :
\begin{enumerate}
  \item \textbf{objectifs} fixés et périmètre ;
  \item \textbf{backlog} planifié (user stories) ;
  \item \textbf{analyse} (conception technique, architecture) ;
  \item \textbf{réalisation} (mise en œuvre, captures d'écran, validation).
\end{enumerate}

\paragraph{Codification des priorités (backlog).}
Dans les tableaux de backlog, la colonne \textbf{Priorité} utilise les codes suivants :
\textbf{H} (haute), \textbf{M} (moyenne) et \textbf{B} (basse).

Le tableau~\ref{tab:r2-overview} synthétise les trois sprints de la release.

\begin{table}[H]
\centering
\caption{Vue d'ensemble de la Release~2}
\label{tab:r2-overview}
\footnotesize
\renewcommand{\arraystretch}{1.15}
\setlength{\tabcolsep}{6pt}
\begin{tabular}{|p{1.2cm}|p{4.2cm}|p{7.8cm}|}
\hline
\textbf{Sprint} & \textbf{Intitulé} & \textbf{Livrables principaux} \\
\hline
Sprint~1 & Mise en place de l'infrastructure &
  Terraform (VNet, NSG, VMs), AKS Engine (cluster K8s), plan d'adressage \\
\hline
Sprint~2 & Intégration continue et déploiement continu (CI/CD) &
  Pipeline Azure DevOps, SonarQube, Docker, Nexus, Trivy, Terraform automatisé \\
\hline
Sprint~3 & Monitoring et supervision de la plateforme &
  Déploiement K8s (CD), MongoDB in-cluster, Prometheus, Grafana (Ansible) \\
\hline
\end{tabular}
\end{table}

\paragraph{Périmètre par rapport à la Release~1.}
La Release~1 couvre le développement fonctionnel (NestJS, Next.js, MongoDB locale,
Cloudinary, notifications WebSocket). La Release~2 ne modifie pas le métier applicatif :
elle prépare et automatise l'environnement cloud, la livraison des images Docker et le
déploiement des composants \texttt{ticket-backend}, \texttt{ticket-frontend} et
\texttt{mongodb} dans le namespace Kubernetes \texttt{ticketmanager}.

\paragraph{Fichiers sources du dépôt.}
\begin{itemize}
  \item Infrastructure : répertoire \texttt{terraform/} ;
  \item Pipeline : \texttt{azure-pipelines.yml} ;
  \item Déploiement Kubernetes : répertoire \texttt{k8s/} ;
  \item Automatisation : répertoire \texttt{ansible/} ;
  \item Documentation Release~1 (socle applicatif) : \texttt{docs/release-1/}.
\end{itemize}

% --- Sprint 1 : Infrastructure ------------------------------------------------
% =============================================================================
% Release 2 — Sprint 1 : Mise en place de l'infrastructure
% =============================================================================

\section[Sprint 1 : Mise en place de l'infrastructure]%
{Sprint 1 : Mise en place de l'infrastructure}

Ce sprint pose les fondations cloud de \textit{TicketManager} sur Azure Stack Hub.
L'objectif est de disposer d'un réseau segmenté, de machines virtuelles d'administration
et de base de données, ainsi que d'un cluster Kubernetes opérationnel, le tout
provisionné de manière reproductible.


% ---------------------------------------------------------------------------
\subsection{Objectif du Sprint 1}
% ---------------------------------------------------------------------------

Le Sprint~1 succède au socle applicatif de la Release~1. Il se concentre sur
\textbf{l'approvisionnement de l'infrastructure} et la \textbf{préparation du cluster}
destiné à héberger l'application conteneurisée. Le tableau~\ref{tab:r2-s1-objectifs} détaille
les objectifs par axe.

\begin{table}[H]
\centering
\caption{Objectifs du Sprint~1 par axe}
\label{tab:r2-s1-objectifs}
\footnotesize
\renewcommand{\arraystretch}{1.1}
\setlength{\tabcolsep}{6pt}
\begin{tabular}{|p{2.8cm}|p{9.2cm}|}
\hline
\textbf{Axe} & \textbf{Objectif} \\
\hline
Infrastructure as Code &
  Décrire et provisionner le réseau, la sécurité et les VMs par Terraform
  (fournisseur \textit{azurestack}) \\
\hline
Segmentation réseau &
  Structurer le VNet \texttt{10.1.0.0/16} en sous-réseaux \texttt{mgmt}, \texttt{db}
  et \texttt{aks}, chacun protégé par un NSG \\
\hline
Machines virtuelles &
  Déployer \textbf{vm-gestion} (administration, agent DevOps) et \textbf{vm-db}
  (MongoDB legacy / migration) \\
\hline
Orchestration &
  Créer le cluster Kubernetes (1 master + 2 workers) via \textbf{AKS Engine} \\
\hline
Vérification &
  Contrôler les ressources dans le portail Azure Stack Hub et via \texttt{kubectl} \\
\hline
\end{tabular}
\end{table}

\paragraph{Résultat attendu.}
À l'issue du sprint :
\begin{itemize}
  \item le réseau virtuel \texttt{ticketmanager-iac-vnet} et les trois sous-réseaux sont
        créés ;
  \item \textbf{vm-gestion} (\texttt{10.1.1.4}) est accessible en SSH et héberge l'agent
        Azure DevOps, Docker, Ansible et AKS Engine ;
  \item \textbf{vm-db} (\texttt{10.1.2.4}) est joignable en IP privée depuis \texttt{mgmt}
        et \texttt{aks} ;
  \item le cluster Kubernetes affiche trois nœuds \texttt{Ready}
        (\texttt{kubectl get nodes}).
\end{itemize}

\paragraph{Périmètre.}
Le tableau~\ref{tab:r2-s1-perimetre} précise le contenu du sprint par rapport aux livraisons
ultérieures.

\begin{table}[H]
\centering
\caption{Périmètre du Sprint~1}
\label{tab:r2-s1-perimetre}
\footnotesize
\renewcommand{\arraystretch}{1.1}
\setlength{\tabcolsep}{6pt}
\begin{tabular}{|p{5.5cm}|c|c|}
\hline
\textbf{Élément} & \textbf{Sprint~1} & \textbf{Sprint ultérieur} \\
\hline
Socle applicatif (Release~1) & Livré & — \\
\hline
Terraform — VNet, NSG, NIC, VMs & Oui & — \\
\hline
AKS Engine — cluster K8s & Oui & — \\
\hline
Pipeline Azure DevOps & — & Sprint~2 \\
\hline
Déploiement pods applicatifs & — & Sprint~3 \\
\hline
Prometheus / Grafana & — & Sprint~3 \\
\hline
\end{tabular}
\end{table}


% ---------------------------------------------------------------------------
\subsection{Backlog du Sprint 1}
% ---------------------------------------------------------------------------

Le backlog du Sprint~1 est organisé en deux epics : infrastructure Terraform et cluster
Kubernetes. Les tableaux~\ref{tab:r2-s1-backlog-infra} et~\ref{tab:r2-s1-backlog-k8s} listent
les user stories planifiées.

\subsubsection{Epic — Infrastructure cloud (Terraform)}

\begin{table}[H]
\centering
\caption{Backlog Sprint~1 — Epic infrastructure Terraform}
\label{tab:r2-s1-backlog-infra}
\footnotesize
\renewcommand{\arraystretch}{0.9}
\setlength{\tabcolsep}{4pt}
\begin{tabular}{|c|p{7.8cm}|c|c|}
\hline
\textbf{ID} & \textbf{User story} & \textbf{Priorité} & \textbf{Est. (j)} \\
\hline
R2-S1-I01 &
  En tant que développeur, je provisionne l'infrastructure par Terraform sur Azure Stack Hub &
  H & 3 \\
\hline
R2-S1-I02 &
  En tant que développeur, je segmente le réseau en trois sous-réseaux dédiés (mgmt, db, aks) &
  H & 2 \\
\hline
R2-S1-I03 &
  En tant que développeur, je configure les groupes de sécurité réseau par sous-réseau &
  H & 2 \\
\hline
R2-S1-I04 &
  En tant que développeur, je dispose d'une VM d'administration (vm-gestion) pour orchestrer les déploiements &
  H & 2 \\
\hline
R2-S1-I05 &
  En tant que développeur, je déploie une VM base de données (vm-db) pour MongoDB Legacy &
  M & 1 \\
\hline
R2-S1-I06 &
  En tant que développeur, j'associe chaque VM à une interface réseau sécurisée &
  H & 1 \\
\hline
\end{tabular}
\end{table}

\subsubsection{Epic — Cluster Kubernetes (AKS Engine)}

\begin{table}[H]
\centering
\caption{Backlog Sprint~1 — Epic cluster Kubernetes}
\label{tab:r2-s1-backlog-k8s}
\footnotesize
\renewcommand{\arraystretch}{0.9}
\setlength{\tabcolsep}{4pt}
\begin{tabular}{|c|p{7.8cm}|c|c|}
\hline
\textbf{ID} & \textbf{User story} & \textbf{Priorité} & \textbf{Est. (j)} \\
\hline
R2-S1-K01 &
  En tant que développeur, je prépare le fichier de configuration AKS Engine (\texttt{kubernetes.json}) &
  H & 2 \\
\hline
R2-S1-K02 &
  En tant que développeur, je déploie le cluster Kubernetes (1 master + 2 workers) &
  H & 3 \\
\hline
R2-S1-K03 &
  En tant que développeur, je vérifie l'état des nœuds du cluster (\texttt{Ready}) &
  H & 1 \\
\hline
R2-S1-K04 &
  En tant que développeur, je configure l'accès \texttt{kubeconfig} depuis vm-gestion &
  H & 1 \\
\hline
\end{tabular}
\end{table}


% ---------------------------------------------------------------------------
\subsection{Conception de l'infrastructure}
% ---------------------------------------------------------------------------

\subsubsection{Plan d'adressage réseau}

Le tableau~\ref{tab:r2-s1-adressage} résume le plan d'adressage retenu pour l'infrastructure
TicketManager. La figure~\ref{fig:r2-s1-plan-adressage} en propose une représentation
schématique (sous-réseaux, adresses IP et séparation Terraform / AKS Engine).

\begin{table}[H]
\centering
\caption{Plan d'adressage réseau TicketManager}
\label{tab:r2-s1-adressage}
\footnotesize
\renewcommand{\arraystretch}{1.1}
\setlength{\tabcolsep}{6pt}
\begin{tabular}{|l|l|}
\hline
\textbf{Élément} & \textbf{Valeur} \\
\hline
Réseau virtuel (VNet) & \texttt{10.1.0.0/16} \\
\hline
Subnet administration (\texttt{mgmt}) & \texttt{10.1.1.0/24} \\
\hline
Subnet base de données (\texttt{db}) & \texttt{10.1.2.0/24} \\
\hline
Subnet Kubernetes (\texttt{aks}) & \texttt{10.1.3.0/24} \\
\hline
IP vm-gestion (privée) & \texttt{10.1.1.4} \\
\hline
IP vm-db (privée) & \texttt{10.1.2.4} \\
\hline
Workers K8s (exemple) & \texttt{10.1.3.4}, \texttt{10.1.3.5} \\
\hline
\end{tabular}
\end{table}

\begin{figure}[htbp]
    \centering
    \includegraphics[width=0.95\textwidth]{plan-adressage-reseau-sprint1.png}
    \caption{Plan d'adressage réseau TicketManager — vue schématique (draw.io)}
    \label{fig:r2-s1-plan-adressage}
\end{figure}

\subsubsection{Architecture globale}

L'architecture repose sur une organisation \textbf{hybride} : couche réseau et compute
provisionnée par Terraform, complétée par un cluster Kubernetes déployé via AKS Engine
sur le sous-réseau \texttt{aks}. La figure~\ref{fig:r2-s1-arch-infra} en donne une vue
d'ensemble.

\begin{figure}[htbp]
    \centering
    \includegraphics[width=0.95\textwidth]{architecture-infrastructure-sprint1.png}
    \caption{Architecture globale de l'infrastructure TicketManager (Sprint~1)}
    \label{fig:r2-s1-arch-infra}
\end{figure}

\paragraph{Composants provisionnés par Terraform.}
\begin{itemize}
  \item \textbf{Réseau virtuel et sous-réseaux} : connectivité interne entre VMs et cluster ;
  \item \textbf{Groupes de sécurité réseau (NSG)} : filtrage SSH, MongoDB, trafic interne ;
  \item \textbf{Interfaces réseau (NIC)} : rattachement des VMs au VNet ;
  \item \textbf{Machines virtuelles} : vm-gestion et vm-db
        (fichiers \texttt{vm-management.tf}, \texttt{vm-db.tf}).
\end{itemize}

\paragraph{Cluster Kubernetes (hors Terraform).}
Terraform prépare le sous-réseau \texttt{aks} et les règles NSG associées ; le déploiement
effectif du plan de contrôle et des workers est réalisé par \textbf{AKS Engine}, outil
open-source adapté à Azure Stack Hub.


% ---------------------------------------------------------------------------
\subsection{Réalisation}
% ---------------------------------------------------------------------------

\subsubsection{Infrastructure as Code avec Terraform}

Terraform adopte une approche \textbf{déclarative} : l'état cible des ressources est décrit
dans des fichiers de configuration ; le moteur calcule et applique les écarts. Le cycle
repose sur quatre commandes : \texttt{terraform init}, \texttt{validate}, \texttt{plan} et
\texttt{apply}.

La figure~\ref{fig:r2-s1-tf-init} présente le résultat de l'initialisation de l'environnement
Terraform. La figure~\ref{fig:r2-s1-tf-validate} confirme la validité syntaxique de la
configuration. La figure~\ref{fig:r2-s1-tf-plan} illustre le plan d'exécution généré avant
application. Enfin, la figure~\ref{fig:r2-s1-tf-apply} montre les sorties du déploiement,
notamment les adresses IP attribuées aux machines virtuelles.

\begin{figure}[htbp]
    \centering
    \includegraphics[width=0.85\textwidth]{s1-terraform-init.png}
    \caption{Initialisation Terraform (\texttt{terraform init})}
    \label{fig:r2-s1-tf-init}
\end{figure}

\begin{figure}[htbp]
    \centering
    \includegraphics[width=0.85\textwidth]{s1-terraform-validate.png}
    \caption{Validation Terraform (\texttt{terraform validate})}
    \label{fig:r2-s1-tf-validate}
\end{figure}

\begin{figure}[htbp]
    \centering
    \includegraphics[width=0.85\textwidth]{s1-terraform-plan.png}
    \caption{Plan d'exécution Terraform (\texttt{terraform plan})}
    \label{fig:r2-s1-tf-plan}
\end{figure}

\begin{figure}[htbp]
    \centering
    \includegraphics[width=0.85\textwidth]{s1-terraform-apply.png}
    \caption{Application Terraform — sorties (IPs vm-gestion, vm-db)}
    \label{fig:r2-s1-tf-apply}
\end{figure}

\paragraph{Vérification portail Azure Stack Hub.}
Les figures~\ref{fig:r2-s1-vnet},~\ref{fig:r2-s1-nsg} et~\ref{fig:r2-s1-vms} confirment la
création des ressources dans le groupe \texttt{AmalBH-RG}. La figure~\ref{fig:r2-s1-vnet}
présente le réseau virtuel et les sous-réseaux. La figure~\ref{fig:r2-s1-nsg} détaille les
groupes de sécurité réseau. La figure~\ref{fig:r2-s1-vms} liste les machines virtuelles
déployées.

\begin{figure}[htbp]
    \centering
    \includegraphics[width=0.85\textwidth]{s1-vnet-subnets.png}
    \caption{Réseau virtuel et sous-réseaux dans le portail}
    \label{fig:r2-s1-vnet}
\end{figure}

\begin{figure}[htbp]
    \centering
    \includegraphics[width=0.85\textwidth]{s1-nsg.png}
    \caption{Groupes de sécurité réseau}
    \label{fig:r2-s1-nsg}
\end{figure}

\begin{figure}[htbp]
    \centering
    \includegraphics[width=0.85\textwidth]{s1-vms-portal.png}
    \caption{Machines virtuelles déployées}
    \label{fig:r2-s1-vms}
\end{figure}


\subsubsection{Déploiement du cluster Kubernetes avec AKS Engine}

\paragraph{Configuration du cluster.}
Le fichier \texttt{kubernetes.json} (modèle API AKS Engine) spécifie la version de
Kubernetes, le plugin réseau (\texttt{kubenet}), le runtime (\texttt{containerd}), la
topologie des nœuds et l'authentification SSH. La figure~\ref{fig:r2-s1-aks-json} présente
un extrait de ce fichier de configuration.

\begin{figure}[htbp]
    \centering
    \includegraphics[width=0.85\textwidth]{s1-aks-engine-json.png}
    \caption{Extrait du fichier de configuration AKS Engine}
    \label{fig:r2-s1-aks-json}
\end{figure}

\paragraph{Exécution du déploiement.}
Le cluster est créé depuis vm-gestion par la commande \texttt{aks-engine deploy}. La
figure~\ref{fig:r2-s1-aks-deploy} illustre l'exécution réussie de cette commande.

\begin{figure}[htbp]
    \centering
    \includegraphics[width=0.85\textwidth]{s1-aks-engine-deploy.png}
    \caption{Déploiement du cluster avec AKS Engine}
    \label{fig:r2-s1-aks-deploy}
\end{figure}

\paragraph{Vérification.}
La figure~\ref{fig:r2-s1-kubectl-nodes} présente la sortie de \texttt{kubectl get nodes} :
les trois nœuds du cluster sont en état \texttt{Ready}.

\begin{figure}[htbp]
    \centering
    \includegraphics[width=0.85\textwidth]{s1-kubectl-nodes.png}
    \caption{Vérification des nœuds Kubernetes (\texttt{kubectl get nodes})}
    \label{fig:r2-s1-kubectl-nodes}
\end{figure}


% ---------------------------------------------------------------------------
\subsection{Bilan du Sprint 1}
% ---------------------------------------------------------------------------

Le Sprint~1 fournit l'infrastructure réseau et compute nécessaire aux sprints suivants :
vm-gestion comme point d'ancrage DevOps, sous-réseau \texttt{aks} prêt pour les pods, et
cluster Kubernetes opérationnel. Le pipeline CI/CD (Sprint~2) et le déploiement applicatif
(Sprint~3) s'appuient sur ce socle.


% --- Sprint 2 : CI/CD ---------------------------------------------------------
% =============================================================================
% Release 2 — Sprint 2 : Intégration continue et déploiement continu (CI/CD)
% =============================================================================

\section[Sprint 2 : Intégration continue et déploiement continu (CI/CD)]%
{Sprint 2 : Intégration continue et déploiement continu (CI/CD)}

Ce sprint industrialise la validation du code, la construction des artefacts conteneurisés
et l'automatisation du provisionnement infrastructurel. Il s'appuie sur \textbf{Azure DevOps}
et un pipeline entièrement décrit en YAML (\texttt{azure-pipelines.yml}).


% ---------------------------------------------------------------------------
\subsection{Objectif du Sprint 2}
% ---------------------------------------------------------------------------

Le tableau~\ref{tab:r2-s2-objectifs} résume les objectifs du Sprint~2 par axe.

\begin{table}[H]
\centering
\caption{Objectifs du Sprint~2 par axe}
\label{tab:r2-s2-objectifs}
\footnotesize
\renewcommand{\arraystretch}{1.1}
\setlength{\tabcolsep}{6pt}
\begin{tabular}{|p{2.8cm}|p{9.2cm}|}
\hline
\textbf{Axe} & \textbf{Objectif} \\
\hline
Intégration continue &
  Pipeline déclenché sur \texttt{main} et \texttt{develop}, exécuté sur agent
  auto-hébergé (\texttt{ticketmanager-pool}) \\
\hline
Qualité &
  Analyse statique SonarQube ; tests automatisés backend (Jest) \\
\hline
Sécurité &
  Scan des images Docker (Trivy) ; analyse des dépendances (Snyk, optionnel) \\
\hline
Conteneurisation &
  Images multi-étapes \texttt{ticket-backend} et \texttt{ticket-frontend} \\
\hline
Registre &
  Publication des images vers Nexus Repository (\texttt{localhost:8082}) \\
\hline
IaC automatisé &
  Exécution Terraform (init, plan, apply) intégrée au pipeline \\
\hline
\end{tabular}
\end{table}

\paragraph{Résultat attendu.}
\begin{itemize}
  \item pipeline Azure DevOps vert de bout en bout (stages Build à Push minimum) ;
  \item images Docker versionnées dans Nexus (\texttt{docker-ticketmanager/}) ;
  \item rapport SonarQube pour le projet \texttt{ticket-manager} ;
  \item infrastructure Terraform synchronisée via le stage IaC.
\end{itemize}

\paragraph{Note sur les stages temporairement désactivés.}
Dans la version actuelle de \texttt{azure-pipelines.yml}, les stages \textbf{Test},
\textbf{Snyk} et \textbf{Trivy} sont commentés pour accélérer les itérations. Le squelette
du rapport prévoit leurs figures ; réactivez-les avant la soutenance si possible.


% ---------------------------------------------------------------------------
\subsection{Backlog du Sprint 2}
% ---------------------------------------------------------------------------

Le tableau~\ref{tab:r2-s2-backlog} présente le backlog planifié pour le Sprint~2.

\subsubsection{Epic — Chaîne d'intégration continue}

\begin{table}[H]
\centering
\caption{Backlog Sprint~2 — Epic intégration continue et déploiement continu}
\label{tab:r2-s2-backlog}
\footnotesize
\renewcommand{\arraystretch}{0.85}
\setlength{\tabcolsep}{4pt}
\begin{tabular}{|c|p{7.8cm}|c|c|}
\hline
\textbf{ID} & \textbf{User story} & \textbf{Priorité} & \textbf{Est. (j)} \\
\hline
R2-S2-C01 &
  En tant que développeur, je centralise le code source sur Azure Repos &
  H & 1 \\
\hline
R2-S2-C02 &
  En tant que développeur, j'exécute le pipeline sur un agent auto-hébergé &
  H & 2 \\
\hline
R2-S2-C03 &
  En tant que développeur, je déclenche automatiquement le pipeline à chaque push &
  H & 1 \\
\hline
R2-S2-C04 &
  En tant que développeur, je compile le backend NestJS et le frontend Next.js à chaque livraison &
  H & 2 \\
\hline
R2-S2-C05 &
  En tant que développeur, j'exécute les tests automatisés et publie la couverture &
  H & 2 \\
\hline
R2-S2-C06 &
  En tant que développeur, j'intègre l'analyse SonarQube dans le pipeline &
  H & 2 \\
\hline
R2-S2-C07 &
  En tant que développeur, je construis des images Docker prêtes pour la production &
  H & 2 \\
\hline
R2-S2-C08 &
  En tant que développeur, je scanne les images Docker avec Trivy avant publication &
  M & 1 \\
\hline
R2-S2-C09 &
  En tant que développeur, je publie les images validées dans Nexus Repository &
  H & 1 \\
\hline
R2-S2-C10 &
  En tant que développeur, j'automatise Terraform dans le pipeline avec validation manuelle &
  H & 2 \\
\hline
\end{tabular}
\end{table}


% ---------------------------------------------------------------------------
\subsection{Analyse — Conception du pipeline}
% ---------------------------------------------------------------------------

\subsubsection{Présentation d'Azure DevOps}

Azure DevOps regroupe les outils collaboratifs mobilisés dans TicketManager :
\begin{itemize}
  \item \textbf{Azure Boards} — planification agile et suivi des sprints ;
  \item \textbf{Azure Repos} — gestion de versions Git ;
  \item \textbf{Azure Pipelines} — exécution automatisée du pipeline CI/CD en YAML.
\end{itemize}

La figure~\ref{fig:r2-s2-ado-ui} présente l'interface principale d'Azure DevOps utilisée
dans le projet.

\begin{figure}[htbp]
    \centering
    \includegraphics[width=0.85\textwidth]{s2-azure-devops-ui.png}
    \caption{Interface principale d'Azure DevOps}
    \label{fig:r2-s2-ado-ui}
\end{figure}

\subsubsection{Structure déclarative du pipeline}

Le pipeline est organisé en \textbf{stages} séquentiels. La figure~\ref{fig:r2-s2-pipeline-stages}
présente la vue d'ensemble de ces étapes telles que définies dans \texttt{azure-pipelines.yml} :

\begin{enumerate}
  \item \textbf{Build} — compilation backend et frontend ;
  \item \textbf{CodeQuality} — analyse SonarQube ;
  \item \textbf{Docker} — construction des images ;
  \item \textbf{Push} — publication vers Nexus ;
  \item \textbf{Terraform} — provisionnement IaC ;
  \item \textbf{DeployK8s} — déploiement Kubernetes (Sprint~3, branche \texttt{main}) ;
  \item \textbf{DeployMonitoring} — Prometheus + Grafana via Ansible (Sprint~3).
\end{enumerate}

\begin{figure}[htbp]
    \centering
    \includegraphics[width=0.95\textwidth]{s2-pipeline-stages.png}
    \caption{Vue d'ensemble des stages du pipeline (\texttt{azure-pipelines.yml})}
    \label{fig:r2-s2-pipeline-stages}
\end{figure}

\subsubsection{Pool d'agents et connexion au dépôt}

Les tâches s'exécutent sur un \textbf{agent auto-hébergé} installé sur vm-gestion, regroupé
dans le pool \texttt{ticketmanager-pool}. La figure~\ref{fig:r2-s2-agent-pool} illustre la
configuration de ce pool d'agents. La figure~\ref{fig:r2-s2-repos} présente la structure du
dépôt Azure Repos, incluant les répertoires \texttt{backend}, \texttt{frontend},
\texttt{terraform} et \texttt{k8s}.

\begin{figure}[htbp]
    \centering
    \includegraphics[width=0.85\textwidth]{s2-agent-pool.png}
    \caption{Configuration du pool d'agents auto-hébergés}
    \label{fig:r2-s2-agent-pool}
\end{figure}

\begin{figure}[htbp]
    \centering
    \includegraphics[width=0.85\textwidth]{s2-azure-repos.png}
    \caption{Structure du dépôt Azure Repos (backend, frontend, terraform, k8s)}
    \label{fig:r2-s2-repos}
\end{figure}


% ---------------------------------------------------------------------------
\subsection{Réalisation}
% ---------------------------------------------------------------------------

\subsubsection{Compilation et tests}

\paragraph{Stage Build.}
Installation des dépendances npm (avec cache local sur vm-gestion), puis
\texttt{npm run build} pour le backend NestJS et le frontend Next.js.

\paragraph{Tests backend.}
Les tests unitaires et d'intégration s'exécutent avec Jest lorsque le stage Test est actif.
La figure~\ref{fig:r2-s2-tests-backend} présente le résultat de l'exécution des tests
automatisés du backend dans Azure DevOps.

\begin{figure}[htbp]
    \centering
    \includegraphics[width=0.85\textwidth]{s2-tests-backend.png}
    \caption{Exécution des tests automatisés du backend}
    \label{fig:r2-s2-tests-backend}
\end{figure}


\subsubsection{Analyse de la qualité du code (SonarQube)}

SonarQube analyse le code source backend (\texttt{backend/src}) et frontend
(\texttt{frontend}). La figure~\ref{fig:r2-s2-sonar-conn} présente la configuration de la
connexion entre SonarQube et Azure DevOps. La figure~\ref{fig:r2-s2-sonar-results} illustre
les résultats d'analyse du projet \texttt{ticket-manager}.

\begin{figure}[htbp]
    \centering
    \includegraphics[width=0.85\textwidth]{s2-sonar-connection.png}
    \caption{Connexion SonarQube — Azure DevOps}
    \label{fig:r2-s2-sonar-conn}
\end{figure}

\begin{figure}[htbp]
    \centering
    \includegraphics[width=0.85\textwidth]{s2-sonar-results.png}
    \caption{Résultats d'analyse SonarQube — projet TicketManager}
    \label{fig:r2-s2-sonar-results}
\end{figure}


\subsubsection{Construction des images Docker}

Deux images sont produites : \textbf{ticket-backend} (API NestJS, port 3001) et
\textbf{ticket-frontend} (Next.js, port 3000). La figure~\ref{fig:r2-s2-dockerfile-be}
présente le Dockerfile multi-étapes du backend. La figure~\ref{fig:r2-s2-dockerfile-fe}
détaille celui du frontend. La figure~\ref{fig:r2-s2-docker-build} montre la construction
réussie des images dans le stage Docker du pipeline.

\begin{figure}[htbp]
    \centering
    \includegraphics[width=0.85\textwidth]{s2-dockerfile-backend.png}
    \caption{Dockerfile du backend NestJS}
    \label{fig:r2-s2-dockerfile-be}
\end{figure}

\begin{figure}[htbp]
    \centering
    \includegraphics[width=0.85\textwidth]{s2-dockerfile-frontend.png}
    \caption{Dockerfile du frontend Next.js}
    \label{fig:r2-s2-dockerfile-fe}
\end{figure}

\begin{figure}[htbp]
    \centering
    \includegraphics[width=0.85\textwidth]{s2-docker-build.png}
    \caption{Construction des images dans le stage Docker}
    \label{fig:r2-s2-docker-build}
\end{figure}


\subsubsection{Scan de vulnérabilités (Trivy)}

Trivy analyse les images construites et signale les CVE par niveau de gravité. La
figure~\ref{fig:r2-s2-trivy-be} présente le rapport de scan pour l'image backend. La
figure~\ref{fig:r2-s2-trivy-fe} illustre celui de l'image frontend.

\begin{figure}[htbp]
    \centering
    \includegraphics[width=0.85\textwidth]{s2-trivy-backend.png}
    \caption{Scan Trivy — image backend}
    \label{fig:r2-s2-trivy-be}
\end{figure}

\begin{figure}[htbp]
    \centering
    \includegraphics[width=0.85\textwidth]{s2-trivy-frontend.png}
    \caption{Scan Trivy — image frontend}
    \label{fig:r2-s2-trivy-fe}
\end{figure}


\subsubsection{Publication vers Nexus Repository}

Nexus joue le rôle de \textbf{registre Docker privé} sur vm-gestion. La
figure~\ref{fig:r2-s2-nexus-repo} présente la création du dépôt Docker hosted. La
figure~\ref{fig:r2-s2-nexus-images} montre les images \texttt{ticket-backend} et
\texttt{ticket-frontend} stockées après exécution du pipeline.

\begin{figure}[htbp]
    \centering
    \includegraphics[width=0.85\textwidth]{s2-nexus-repo.png}
    \caption{Création du dépôt Docker hosted dans Nexus}
    \label{fig:r2-s2-nexus-repo}
\end{figure}

\begin{figure}[htbp]
    \centering
    \includegraphics[width=0.85\textwidth]{s2-nexus-images.png}
    \caption{Images \texttt{ticket-backend} et \texttt{ticket-frontend} dans Nexus}
    \label{fig:r2-s2-nexus-images}
\end{figure}


\subsubsection{Terraform automatisé dans le pipeline}

Le stage \textbf{Terraform IaC} exécute \texttt{terraform init}, \texttt{plan} et
\texttt{apply}. La figure~\ref{fig:r2-s2-tf-pipeline} présente l'exécution réussie de ce
stage dans Azure DevOps. La figure~\ref{fig:r2-s2-pipeline-run} illustre une exécution
complète du pipeline CI (stages Build à Push).

\begin{figure}[htbp]
    \centering
    \includegraphics[width=0.85\textwidth]{s2-terraform-pipeline.png}
    \caption{Exécution réussie du stage Terraform dans Azure DevOps}
    \label{fig:r2-s2-tf-pipeline}
\end{figure}

\begin{figure}[htbp]
    \centering
    \includegraphics[width=0.85\textwidth]{s2-pipeline-run-success.png}
    \caption{Exécution complète du pipeline CI (stages Build à Push)}
    \label{fig:r2-s2-pipeline-run}
\end{figure}


% ---------------------------------------------------------------------------
\subsection{Bilan du Sprint 2}
% ---------------------------------------------------------------------------

Le Sprint~2 met en place une chaîne d'intégration continue couvrant la compilation, la
qualité, la conteneurisation, la sécurité des images et la publication vers un registre
privé. Le déploiement effectif des pods applicatifs sur le cluster (stage
\textbf{DeployK8s}) et la supervision (stage \textbf{DeployMonitoring}) sont détaillés au
Sprint~3.


% --- Sprint 3 : Monitoring et CD ----------------------------------------------
% =============================================================================
% Release 2 — Sprint 3 : Monitoring et supervision de la plateforme
% =============================================================================

\section[Sprint 3 : Monitoring et supervision de la plateforme]%
{Sprint 3 : Monitoring et supervision de la plateforme}

Ce sprint finalise la Release~2 par le \textbf{déploiement continu} (CD) de l'application
sur Kubernetes et la mise en place de la \textbf{supervision} (Prometheus, Grafana) via
Ansible sur vm-gestion.


% ---------------------------------------------------------------------------
\subsection{Objectif du Sprint 3}
% ---------------------------------------------------------------------------

Le tableau~\ref{tab:r2-s3-objectifs} détaille les objectifs du Sprint~3 par axe.

\begin{table}[H]
\centering
\caption{Objectifs du Sprint~3 par axe}
\label{tab:r2-s3-objectifs}
\footnotesize
\renewcommand{\arraystretch}{1.1}
\setlength{\tabcolsep}{6pt}
\begin{tabular}{|p{2.8cm}|p{9.2cm}|}
\hline
\textbf{Axe} & \textbf{Objectif} \\
\hline
Déploiement continu &
  Déployer automatiquement backend, frontend et MongoDB in-cluster via le stage
  \texttt{DeployK8s} (branche \texttt{main}, approbation manuelle) \\
\hline
Manifests Kubernetes &
  Appliquer les fichiers YAML du répertoire \texttt{k8s/} (Deployments, Services,
  StatefulSet MongoDB, Secrets) \\
\hline
Registre privé &
  Configurer le pull d'images depuis Nexus (\texttt{10.1.1.4:8082}) sur les workers \\
\hline
Supervision &
  Déployer Prometheus et Grafana sur vm-gestion (\texttt{ansible/playbooks/monitoring.yml}) \\
\hline
Observabilité &
  Collecter les métriques des nœuds K8s, du backend NestJS et de MongoDB \\
\hline
Validation &
  Smoke tests post-déploiement (API, frontend, cibles Prometheus UP) \\
\hline
\end{tabular}
\end{table}

\paragraph{Résultat attendu.}
\begin{itemize}
  \item namespace \texttt{ticketmanager} avec pods \texttt{Running} ;
  \item application accessible via NodePort (frontend \texttt{30080}, backend \texttt{30081}) ;
  \item Prometheus accessible sur le port \texttt{9090}, Grafana sur \texttt{3030} ;
  \item dashboards opérationnels pour le cluster et l'application.
\end{itemize}

\paragraph{Choix architectural — MongoDB in-cluster.}
Contrairement aux projets précédents (chatbot) où MongoDB était exclusivement sur une VM
dédiée, TicketManager héberge la base de données de production dans un
\textbf{StatefulSet Kubernetes} avec volume persistant. La vm-db reste disponible pour
migration ou sauvegarde legacy (\texttt{ansible/playbooks/mongodb.yml}).


% ---------------------------------------------------------------------------
\subsection{Backlog du Sprint 3}
% ---------------------------------------------------------------------------

Le tableau~\ref{tab:r2-s3-backlog-deploy} couvre l'epic déploiement continu. Le
tableau~\ref{tab:r2-s3-backlog-monitoring} couvre l'epic supervision.

\subsubsection{Epic — Déploiement continu (Kubernetes)}

\begin{table}[H]
\centering
\caption{Backlog Sprint~3 — Epic déploiement continu Kubernetes}
\label{tab:r2-s3-backlog-deploy}
\footnotesize
\renewcommand{\arraystretch}{0.85}
\setlength{\tabcolsep}{4pt}
\begin{tabular}{|c|p{7.8cm}|c|c|}
\hline
\textbf{ID} & \textbf{User story} & \textbf{Priorité} & \textbf{Est. (j)} \\
\hline
R2-S3-D01 &
  En tant que développeur, je déploie le backend NestJS en pod Kubernetes &
  H & 2 \\
\hline
R2-S3-D02 &
  En tant que développeur, je déploie le frontend Next.js en pod Kubernetes &
  H & 2 \\
\hline
R2-S3-D03 &
  En tant que développeur, je déploie MongoDB in-cluster via un StatefulSet &
  H & 2 \\
\hline
R2-S3-D04 &
  En tant que développeur, je configure les secrets Kubernetes (JWT, MongoDB, Cloudinary) &
  H & 1 \\
\hline
R2-S3-D05 &
  En tant que développeur, j'automatise le déploiement via le stage CD du pipeline &
  H & 2 \\
\hline
R2-S3-D06 &
  En tant que développeur, je valide le déploiement par des smoke tests automatisés &
  H & 1 \\
\hline
\end{tabular}
\end{table}

\subsubsection{Epic — Monitoring et supervision}

\begin{table}[H]
\centering
\caption{Backlog Sprint~3 — Epic monitoring et supervision}
\label{tab:r2-s3-backlog-monitoring}
\footnotesize
\renewcommand{\arraystretch}{0.85}
\setlength{\tabcolsep}{4pt}
\begin{tabular}{|c|p{7.8cm}|c|c|}
\hline
\textbf{ID} & \textbf{User story} & \textbf{Priorité} & \textbf{Est. (j)} \\
\hline
R2-S3-M01 &
  En tant que développeur, je déploie Prometheus sur vm-gestion via Ansible &
  H & 2 \\
\hline
R2-S3-M02 &
  En tant que développeur, je déploie Grafana et je connecte Prometheus comme source de données &
  H & 2 \\
\hline
R2-S3-M03 &
  En tant que développeur, j'intègre le déploiement monitoring au pipeline Azure DevOps &
  M & 1 \\
\hline
R2-S3-M04 &
  En tant que développeur, je vérifie que les cibles Prometheus sont en état \texttt{UP} &
  H & 1 \\
\hline
\end{tabular}
\end{table}


% ---------------------------------------------------------------------------
\subsection{Analyse — Conception du déploiement Kubernetes}
% ---------------------------------------------------------------------------

\subsubsection{Architecture de déploiement}

Les composants applicatifs sont encapsulés dans des pods du namespace
\texttt{ticketmanager} :
\begin{itemize}
  \item \textbf{ticket-backend} — image Nexus, port 3001, secrets JWT et Cloudinary ;
  \item \textbf{ticket-frontend} — image Nexus, port 3000, URL API figée au build Docker ;
  \item \textbf{mongodb} — StatefulSet avec volume persistant pour les données tickets.
\end{itemize}

Les requêtes externes atteignent les services via \textbf{NodePort} sur les workers
(\texttt{30080} frontend, \texttt{30081} backend). La figure~\ref{fig:r2-s3-k8s-arch}
détaille l'architecture de déploiement Kubernetes (ressources du namespace
\texttt{ticketmanager}), en complément de la vue d'ensemble physique (figure~2.5).

\begin{figure}[htbp]
    \centering
    \includegraphics[width=0.92\textwidth]{architecture-deploiement-k8s.png}
    \caption{Architecture de déploiement Kubernetes — namespace \texttt{ticketmanager}}
    \label{fig:r2-s3-k8s-arch}
\end{figure}

\subsubsection{Chaîne de déploiement continu}

Le stage \texttt{DeployK8s} du pipeline :
\begin{enumerate}
  \item attend une \textbf{approbation manuelle} (\texttt{ManualValidation}) ;
  \item configure le registre insecure sur les workers ;
  \item exécute \texttt{k8s/deploy-pipeline.sh} (secrets, manifests, rollout, seed, smoke test).
\end{enumerate}


% ---------------------------------------------------------------------------
\subsection{Réalisation — Déploiement Kubernetes}
% ---------------------------------------------------------------------------

\subsubsection{Manifests et secrets}

La figure~\ref{fig:r2-s3-backend-yaml} présente le manifest de déploiement du backend
(\texttt{backend-deployment.yaml}). La figure~\ref{fig:r2-s3-frontend-yaml} illustre celui
du frontend (\texttt{frontend-deployment.yaml}).

\begin{figure}[htbp]
    \centering
    \includegraphics[width=0.85\textwidth]{s3-backend-deployment.png}
    \caption{Manifest de déploiement du backend (\texttt{backend-deployment.yaml})}
    \label{fig:r2-s3-backend-yaml}
\end{figure}

\begin{figure}[htbp]
    \centering
    \includegraphics[width=0.85\textwidth]{s3-frontend-deployment.png}
    \caption{Manifest de déploiement du frontend (\texttt{frontend-deployment.yaml})}
    \label{fig:r2-s3-frontend-yaml}
\end{figure}

\subsubsection{Exécution du déploiement}

La figure~\ref{fig:r2-s3-k8s-pipeline} présente le stage \textbf{DeployK8s} exécuté dans
Azure DevOps, incluant l'approbation manuelle et l'application des manifests.

\begin{figure}[htbp]
    \centering
    \includegraphics[width=0.85\textwidth]{s3-k8s-pipeline-stage.png}
    \caption{Stage DeployK8s dans Azure DevOps}
    \label{fig:r2-s3-k8s-pipeline}
\end{figure}

\subsubsection{Vérification des pods}

La figure~\ref{fig:r2-s3-pods} présente la sortie de \texttt{kubectl get pods -n ticketmanager}
: l'ensemble des pods applicatifs est en état \texttt{Running}.

\begin{figure}[htbp]
    \centering
    \includegraphics[width=0.85\textwidth]{s3-kubectl-pods.png}
    \caption{État des pods (\texttt{kubectl get pods -n ticketmanager})}
    \label{fig:r2-s3-pods}
\end{figure}

\subsubsection{Application accessible}

La figure~\ref{fig:r2-s3-app} illustre l'interface TicketManager accessible depuis un
navigateur après déploiement réussi sur le cluster.

\begin{figure}[htbp]
    \centering
    \includegraphics[width=0.85\textwidth]{s3-app-browser.png}
    \caption{Interface TicketManager accessible après déploiement}
    \label{fig:r2-s3-app}
\end{figure}


% ---------------------------------------------------------------------------
\subsection{Réalisation — Monitoring et supervision}
% ---------------------------------------------------------------------------

\subsubsection{Déploiement Ansible (Prometheus + Grafana)}

Le playbook \texttt{ansible/playbooks/monitoring.yml} déploie Prometheus (port 9090) et
Grafana (port 3030) sur vm-gestion. La figure~\ref{fig:r2-s3-ansible} présente l'exécution
réussie de ce playbook.

\begin{figure}[htbp]
    \centering
    \includegraphics[width=0.85\textwidth]{s3-ansible-monitoring.png}
    \caption{Exécution du playbook \texttt{monitoring.yml}}
    \label{fig:r2-s3-ansible}
\end{figure}

\subsubsection{Configuration Prometheus}

Le fichier \texttt{ansible/templates/prometheus.yml.j2} définit les cibles surveillées.
La figure~\ref{fig:r2-s3-prom-targets} présente l'interface Prometheus (\texttt{/targets})
avec l'ensemble des cibles en état \texttt{UP}.

\begin{figure}[htbp]
    \centering
    \includegraphics[width=0.85\textwidth]{s3-prometheus-targets.png}
    \caption{Cibles actives dans l'interface Prometheus (\texttt{/targets})}
    \label{fig:r2-s3-prom-targets}
\end{figure}

\subsubsection{Dashboards Grafana}

La figure~\ref{fig:r2-s3-grafana-ds} présente la configuration de Prometheus comme source
de données dans Grafana. La figure~\ref{fig:r2-s3-grafana-dash} illustre un dashboard de
supervision regroupant les métriques du cluster et de l'application backend.

\begin{figure}[htbp]
    \centering
    \includegraphics[width=0.85\textwidth]{s3-grafana-datasource.png}
    \caption{Configuration de Prometheus comme source Grafana}
    \label{fig:r2-s3-grafana-ds}
\end{figure}

\begin{figure}[htbp]
    \centering
    \includegraphics[width=0.85\textwidth]{s3-grafana-dashboard.png}
    \caption{Dashboard de supervision — cluster et application}
    \label{fig:r2-s3-grafana-dash}
\end{figure}

\subsubsection{Intégration pipeline}

La figure~\ref{fig:r2-s3-mon-pipeline} présente le stage \textbf{DeployMonitoring} du
pipeline Azure DevOps, exécutant automatiquement le playbook Ansible après un déploiement
K8s réussi.

\begin{figure}[htbp]
    \centering
    \includegraphics[width=0.85\textwidth]{s3-monitoring-pipeline.png}
    \caption{Stage DeployMonitoring dans Azure DevOps}
    \label{fig:r2-s3-mon-pipeline}
\end{figure}


% ---------------------------------------------------------------------------
\subsection{Tests et validation}
% ---------------------------------------------------------------------------

\subsubsection{Smoke tests post-déploiement}

Le script \texttt{k8s/post-deploy.sh} vérifie la disponibilité de l'API, du frontend et
des outils de supervision. Le tableau~\ref{tab:r2-s3-validation} résume les critères de
validation retenus pour le Sprint~3.

\begin{table}[H]
\centering
\caption{Critères de validation du Sprint~3}
\label{tab:r2-s3-validation}
\footnotesize
\renewcommand{\arraystretch}{1.1}
\setlength{\tabcolsep}{6pt}
\begin{tabular}{|l|l|l|}
\hline
\textbf{Composant} & \textbf{URL / commande} & \textbf{Critère de succès} \\
\hline
Backend & \texttt{curl http://<worker>:30081/} & HTTP 200 \\
\hline
Frontend & \texttt{http://<worker>:30080} & Page de connexion affichée \\
\hline
Pods & \texttt{kubectl get pods -n ticketmanager} & Tous \texttt{Running} \\
\hline
Prometheus & \texttt{curl http://10.1.1.4:9090/-/ready} & \texttt{Ready} \\
\hline
Grafana & \texttt{http://10.1.1.4:3030/login} & HTTP 200 \\
\hline
\end{tabular}
\end{table}


% ---------------------------------------------------------------------------
\subsection{Bilan du Sprint 3}
% ---------------------------------------------------------------------------

Le Sprint~3 clôt la Release~2 en rendant TicketManager déployable et observable en
environnement cloud. L'application tourne dans Kubernetes ; la supervision Prometheus/Grafana
permet de détecter les anomalies de charge, de disponibilité et de performance avant
intervention manuelle.


% --- Bilan Release 2 ----------------------------------------------------------
\subsection{Bilan de la Release~2}

À l'issue des trois sprints, la plateforme \textit{TicketManager} dispose d'une
infrastructure documentée et reproductible sur Azure Stack Hub, d'une chaîne CI/CD
automatisée (compilation, qualité, conteneurisation, publication Nexus, provisionnement
IaC) et d'un déploiement applicatif supervisé (pods Kubernetes, Prometheus, Grafana).

Les indicateurs de réussite retenus pour cette release sont :
\begin{itemize}
  \item exécution réussie du pipeline Azure DevOps sur la branche \texttt{main} ;
  \item pods du namespace \texttt{ticketmanager} en état \texttt{Running} ;
  \item accessibilité de l'interface web et de l'API depuis le réseau autorisé ;
  \item cibles Prometheus en état \texttt{UP} et dashboards Grafana opérationnels.
\end{itemize}

Les évolutions ultérieures (hors Release~2) peuvent porter sur le renforcement des tests
automatisés dans le pipeline, la réactivation des stages Snyk/Trivy, l'exposition via
Ingress Controller et la haute disponibilité multi-réplicas.

%
% Convention : les \caption des tableaux sont placés AU-DESSUS du tabular.
% Les figures de réalisation sont citées dans le texte via \ref{fig:...}.
%
% =============================================================================

\section*{Release 2 : Industrialisation et déploiement de la plateforme}
\addcontentsline{toc}{section}{Release 2 : Industrialisation et déploiement de la plateforme}

La \textbf{Release~2} prolonge le socle applicatif livré en Release~1 (conception et
développement de \textit{TicketManager}). Elle vise à industrialiser le cycle de vie de la
plateforme sur \textbf{Azure Stack Hub} : provisionnement reproductible de l'infrastructure,
chaîne d'intégration et de déploiement continus (CI/CD), déploiement conteneurisé sur
Kubernetes et supervision de l'environnement de production.

\paragraph{Organisation en trois sprints.}
Chaque sprint suit la même trame pédagogique que la Release~1 :
\begin{enumerate}
  \item \textbf{objectifs} fixés et périmètre ;
  \item \textbf{backlog} planifié (user stories) ;
  \item \textbf{analyse} (conception technique, architecture) ;
  \item \textbf{réalisation} (mise en œuvre, captures d'écran, validation).
\end{enumerate}

\paragraph{Codification des priorités (backlog).}
Dans les tableaux de backlog, la colonne \textbf{Priorité} utilise les codes suivants :
\textbf{H} (haute), \textbf{M} (moyenne) et \textbf{B} (basse).

Le tableau~\ref{tab:r2-overview} synthétise les trois sprints de la release.

\begin{table}[H]
\centering
\caption{Vue d'ensemble de la Release~2}
\label{tab:r2-overview}
\footnotesize
\renewcommand{\arraystretch}{1.15}
\setlength{\tabcolsep}{6pt}
\begin{tabular}{|p{1.2cm}|p{4.2cm}|p{7.8cm}|}
\hline
\textbf{Sprint} & \textbf{Intitulé} & \textbf{Livrables principaux} \\
\hline
Sprint~1 & Mise en place de l'infrastructure &
  Terraform (VNet, NSG, VMs), AKS Engine (cluster K8s), plan d'adressage \\
\hline
Sprint~2 & Intégration continue et déploiement continu (CI/CD) &
  Pipeline Azure DevOps, SonarQube, Docker, Nexus, Trivy, Terraform automatisé \\
\hline
Sprint~3 & Monitoring et supervision de la plateforme &
  Déploiement K8s (CD), MongoDB in-cluster, Prometheus, Grafana (Ansible) \\
\hline
\end{tabular}
\end{table}

\paragraph{Périmètre par rapport à la Release~1.}
La Release~1 couvre le développement fonctionnel (NestJS, Next.js, MongoDB locale,
Cloudinary, notifications WebSocket). La Release~2 ne modifie pas le métier applicatif :
elle prépare et automatise l'environnement cloud, la livraison des images Docker et le
déploiement des composants \texttt{ticket-backend}, \texttt{ticket-frontend} et
\texttt{mongodb} dans le namespace Kubernetes \texttt{ticketmanager}.

\paragraph{Fichiers sources du dépôt.}
\begin{itemize}
  \item Infrastructure : répertoire \texttt{terraform/} ;
  \item Pipeline : \texttt{azure-pipelines.yml} ;
  \item Déploiement Kubernetes : répertoire \texttt{k8s/} ;
  \item Automatisation : répertoire \texttt{ansible/} ;
  \item Documentation Release~1 (socle applicatif) : \texttt{docs/release-1/}.
\end{itemize}

% --- Sprint 1 : Infrastructure ------------------------------------------------
% =============================================================================
% Release 2 — Sprint 1 : Mise en place de l'infrastructure
% =============================================================================

\section[Sprint 1 : Mise en place de l'infrastructure]%
{Sprint 1 : Mise en place de l'infrastructure}

Ce sprint pose les fondations cloud de \textit{TicketManager} sur Azure Stack Hub.
L'objectif est de disposer d'un réseau segmenté, de machines virtuelles d'administration
et de base de données, ainsi que d'un cluster Kubernetes opérationnel, le tout
provisionné de manière reproductible.


% ---------------------------------------------------------------------------
\subsection{Objectif du Sprint 1}
% ---------------------------------------------------------------------------

Le Sprint~1 succède au socle applicatif de la Release~1. Il se concentre sur
\textbf{l'approvisionnement de l'infrastructure} et la \textbf{préparation du cluster}
destiné à héberger l'application conteneurisée. Le tableau~\ref{tab:r2-s1-objectifs} détaille
les objectifs par axe.

\begin{table}[H]
\centering
\caption{Objectifs du Sprint~1 par axe}
\label{tab:r2-s1-objectifs}
\footnotesize
\renewcommand{\arraystretch}{1.1}
\setlength{\tabcolsep}{6pt}
\begin{tabular}{|p{2.8cm}|p{9.2cm}|}
\hline
\textbf{Axe} & \textbf{Objectif} \\
\hline
Infrastructure as Code &
  Décrire et provisionner le réseau, la sécurité et les VMs par Terraform
  (fournisseur \textit{azurestack}) \\
\hline
Segmentation réseau &
  Structurer le VNet \texttt{10.1.0.0/16} en sous-réseaux \texttt{mgmt}, \texttt{db}
  et \texttt{aks}, chacun protégé par un NSG \\
\hline
Machines virtuelles &
  Déployer \textbf{vm-gestion} (administration, agent DevOps) et \textbf{vm-db}
  (MongoDB legacy / migration) \\
\hline
Orchestration &
  Créer le cluster Kubernetes (1 master + 2 workers) via \textbf{AKS Engine} \\
\hline
Vérification &
  Contrôler les ressources dans le portail Azure Stack Hub et via \texttt{kubectl} \\
\hline
\end{tabular}
\end{table}

\paragraph{Résultat attendu.}
À l'issue du sprint :
\begin{itemize}
  \item le réseau virtuel \texttt{ticketmanager-iac-vnet} et les trois sous-réseaux sont
        créés ;
  \item \textbf{vm-gestion} (\texttt{10.1.1.4}) est accessible en SSH et héberge l'agent
        Azure DevOps, Docker, Ansible et AKS Engine ;
  \item \textbf{vm-db} (\texttt{10.1.2.4}) est joignable en IP privée depuis \texttt{mgmt}
        et \texttt{aks} ;
  \item le cluster Kubernetes affiche trois nœuds \texttt{Ready}
        (\texttt{kubectl get nodes}).
\end{itemize}

\paragraph{Périmètre.}
Le tableau~\ref{tab:r2-s1-perimetre} précise le contenu du sprint par rapport aux livraisons
ultérieures.

\begin{table}[H]
\centering
\caption{Périmètre du Sprint~1}
\label{tab:r2-s1-perimetre}
\footnotesize
\renewcommand{\arraystretch}{1.1}
\setlength{\tabcolsep}{6pt}
\begin{tabular}{|p{5.5cm}|c|c|}
\hline
\textbf{Élément} & \textbf{Sprint~1} & \textbf{Sprint ultérieur} \\
\hline
Socle applicatif (Release~1) & Livré & — \\
\hline
Terraform — VNet, NSG, NIC, VMs & Oui & — \\
\hline
AKS Engine — cluster K8s & Oui & — \\
\hline
Pipeline Azure DevOps & — & Sprint~2 \\
\hline
Déploiement pods applicatifs & — & Sprint~3 \\
\hline
Prometheus / Grafana & — & Sprint~3 \\
\hline
\end{tabular}
\end{table}


% ---------------------------------------------------------------------------
\subsection{Backlog du Sprint 1}
% ---------------------------------------------------------------------------

Le backlog du Sprint~1 est organisé en deux epics : infrastructure Terraform et cluster
Kubernetes. Les tableaux~\ref{tab:r2-s1-backlog-infra} et~\ref{tab:r2-s1-backlog-k8s} listent
les user stories planifiées.

\subsubsection{Epic — Infrastructure cloud (Terraform)}

\begin{table}[H]
\centering
\caption{Backlog Sprint~1 — Epic infrastructure Terraform}
\label{tab:r2-s1-backlog-infra}
\footnotesize
\renewcommand{\arraystretch}{0.9}
\setlength{\tabcolsep}{4pt}
\begin{tabular}{|c|p{7.8cm}|c|c|}
\hline
\textbf{ID} & \textbf{User story} & \textbf{Priorité} & \textbf{Est. (j)} \\
\hline
R2-S1-I01 &
  En tant que développeur, je provisionne l'infrastructure par Terraform sur Azure Stack Hub &
  H & 3 \\
\hline
R2-S1-I02 &
  En tant que développeur, je segmente le réseau en trois sous-réseaux dédiés (mgmt, db, aks) &
  H & 2 \\
\hline
R2-S1-I03 &
  En tant que développeur, je configure les groupes de sécurité réseau par sous-réseau &
  H & 2 \\
\hline
R2-S1-I04 &
  En tant que développeur, je dispose d'une VM d'administration (vm-gestion) pour orchestrer les déploiements &
  H & 2 \\
\hline
R2-S1-I05 &
  En tant que développeur, je déploie une VM base de données (vm-db) pour MongoDB Legacy &
  M & 1 \\
\hline
R2-S1-I06 &
  En tant que développeur, j'associe chaque VM à une interface réseau sécurisée &
  H & 1 \\
\hline
\end{tabular}
\end{table}

\subsubsection{Epic — Cluster Kubernetes (AKS Engine)}

\begin{table}[H]
\centering
\caption{Backlog Sprint~1 — Epic cluster Kubernetes}
\label{tab:r2-s1-backlog-k8s}
\footnotesize
\renewcommand{\arraystretch}{0.9}
\setlength{\tabcolsep}{4pt}
\begin{tabular}{|c|p{7.8cm}|c|c|}
\hline
\textbf{ID} & \textbf{User story} & \textbf{Priorité} & \textbf{Est. (j)} \\
\hline
R2-S1-K01 &
  En tant que développeur, je prépare le fichier de configuration AKS Engine (\texttt{kubernetes.json}) &
  H & 2 \\
\hline
R2-S1-K02 &
  En tant que développeur, je déploie le cluster Kubernetes (1 master + 2 workers) &
  H & 3 \\
\hline
R2-S1-K03 &
  En tant que développeur, je vérifie l'état des nœuds du cluster (\texttt{Ready}) &
  H & 1 \\
\hline
R2-S1-K04 &
  En tant que développeur, je configure l'accès \texttt{kubeconfig} depuis vm-gestion &
  H & 1 \\
\hline
\end{tabular}
\end{table}


% ---------------------------------------------------------------------------
\subsection{Conception de l'infrastructure}
% ---------------------------------------------------------------------------

\subsubsection{Plan d'adressage réseau}

Le tableau~\ref{tab:r2-s1-adressage} résume le plan d'adressage retenu pour l'infrastructure
TicketManager. La figure~\ref{fig:r2-s1-plan-adressage} en propose une représentation
schématique (sous-réseaux, adresses IP et séparation Terraform / AKS Engine).

\begin{table}[H]
\centering
\caption{Plan d'adressage réseau TicketManager}
\label{tab:r2-s1-adressage}
\footnotesize
\renewcommand{\arraystretch}{1.1}
\setlength{\tabcolsep}{6pt}
\begin{tabular}{|l|l|}
\hline
\textbf{Élément} & \textbf{Valeur} \\
\hline
Réseau virtuel (VNet) & \texttt{10.1.0.0/16} \\
\hline
Subnet administration (\texttt{mgmt}) & \texttt{10.1.1.0/24} \\
\hline
Subnet base de données (\texttt{db}) & \texttt{10.1.2.0/24} \\
\hline
Subnet Kubernetes (\texttt{aks}) & \texttt{10.1.3.0/24} \\
\hline
IP vm-gestion (privée) & \texttt{10.1.1.4} \\
\hline
IP vm-db (privée) & \texttt{10.1.2.4} \\
\hline
Workers K8s (exemple) & \texttt{10.1.3.4}, \texttt{10.1.3.5} \\
\hline
\end{tabular}
\end{table}

\begin{figure}[htbp]
    \centering
    \includegraphics[width=0.95\textwidth]{plan-adressage-reseau-sprint1.png}
    \caption{Plan d'adressage réseau TicketManager — vue schématique (draw.io)}
    \label{fig:r2-s1-plan-adressage}
\end{figure}

\subsubsection{Architecture globale}

L'architecture repose sur une organisation \textbf{hybride} : couche réseau et compute
provisionnée par Terraform, complétée par un cluster Kubernetes déployé via AKS Engine
sur le sous-réseau \texttt{aks}. La figure~\ref{fig:r2-s1-arch-infra} en donne une vue
d'ensemble.

\begin{figure}[htbp]
    \centering
    \includegraphics[width=0.95\textwidth]{architecture-infrastructure-sprint1.png}
    \caption{Architecture globale de l'infrastructure TicketManager (Sprint~1)}
    \label{fig:r2-s1-arch-infra}
\end{figure}

\paragraph{Composants provisionnés par Terraform.}
\begin{itemize}
  \item \textbf{Réseau virtuel et sous-réseaux} : connectivité interne entre VMs et cluster ;
  \item \textbf{Groupes de sécurité réseau (NSG)} : filtrage SSH, MongoDB, trafic interne ;
  \item \textbf{Interfaces réseau (NIC)} : rattachement des VMs au VNet ;
  \item \textbf{Machines virtuelles} : vm-gestion et vm-db
        (fichiers \texttt{vm-management.tf}, \texttt{vm-db.tf}).
\end{itemize}

\paragraph{Cluster Kubernetes (hors Terraform).}
Terraform prépare le sous-réseau \texttt{aks} et les règles NSG associées ; le déploiement
effectif du plan de contrôle et des workers est réalisé par \textbf{AKS Engine}, outil
open-source adapté à Azure Stack Hub.


% ---------------------------------------------------------------------------
\subsection{Réalisation}
% ---------------------------------------------------------------------------

\subsubsection{Infrastructure as Code avec Terraform}

Terraform adopte une approche \textbf{déclarative} : l'état cible des ressources est décrit
dans des fichiers de configuration ; le moteur calcule et applique les écarts. Le cycle
repose sur quatre commandes : \texttt{terraform init}, \texttt{validate}, \texttt{plan} et
\texttt{apply}.

La figure~\ref{fig:r2-s1-tf-init} présente le résultat de l'initialisation de l'environnement
Terraform. La figure~\ref{fig:r2-s1-tf-validate} confirme la validité syntaxique de la
configuration. La figure~\ref{fig:r2-s1-tf-plan} illustre le plan d'exécution généré avant
application. Enfin, la figure~\ref{fig:r2-s1-tf-apply} montre les sorties du déploiement,
notamment les adresses IP attribuées aux machines virtuelles.

\begin{figure}[htbp]
    \centering
    \includegraphics[width=0.85\textwidth]{s1-terraform-init.png}
    \caption{Initialisation Terraform (\texttt{terraform init})}
    \label{fig:r2-s1-tf-init}
\end{figure}

\begin{figure}[htbp]
    \centering
    \includegraphics[width=0.85\textwidth]{s1-terraform-validate.png}
    \caption{Validation Terraform (\texttt{terraform validate})}
    \label{fig:r2-s1-tf-validate}
\end{figure}

\begin{figure}[htbp]
    \centering
    \includegraphics[width=0.85\textwidth]{s1-terraform-plan.png}
    \caption{Plan d'exécution Terraform (\texttt{terraform plan})}
    \label{fig:r2-s1-tf-plan}
\end{figure}

\begin{figure}[htbp]
    \centering
    \includegraphics[width=0.85\textwidth]{s1-terraform-apply.png}
    \caption{Application Terraform — sorties (IPs vm-gestion, vm-db)}
    \label{fig:r2-s1-tf-apply}
\end{figure}

\paragraph{Vérification portail Azure Stack Hub.}
Les figures~\ref{fig:r2-s1-vnet},~\ref{fig:r2-s1-nsg} et~\ref{fig:r2-s1-vms} confirment la
création des ressources dans le groupe \texttt{AmalBH-RG}. La figure~\ref{fig:r2-s1-vnet}
présente le réseau virtuel et les sous-réseaux. La figure~\ref{fig:r2-s1-nsg} détaille les
groupes de sécurité réseau. La figure~\ref{fig:r2-s1-vms} liste les machines virtuelles
déployées.

\begin{figure}[htbp]
    \centering
    \includegraphics[width=0.85\textwidth]{s1-vnet-subnets.png}
    \caption{Réseau virtuel et sous-réseaux dans le portail}
    \label{fig:r2-s1-vnet}
\end{figure}

\begin{figure}[htbp]
    \centering
    \includegraphics[width=0.85\textwidth]{s1-nsg.png}
    \caption{Groupes de sécurité réseau}
    \label{fig:r2-s1-nsg}
\end{figure}

\begin{figure}[htbp]
    \centering
    \includegraphics[width=0.85\textwidth]{s1-vms-portal.png}
    \caption{Machines virtuelles déployées}
    \label{fig:r2-s1-vms}
\end{figure}


\subsubsection{Déploiement du cluster Kubernetes avec AKS Engine}

\paragraph{Configuration du cluster.}
Le fichier \texttt{kubernetes.json} (modèle API AKS Engine) spécifie la version de
Kubernetes, le plugin réseau (\texttt{kubenet}), le runtime (\texttt{containerd}), la
topologie des nœuds et l'authentification SSH. La figure~\ref{fig:r2-s1-aks-json} présente
un extrait de ce fichier de configuration.

\begin{figure}[htbp]
    \centering
    \includegraphics[width=0.85\textwidth]{s1-aks-engine-json.png}
    \caption{Extrait du fichier de configuration AKS Engine}
    \label{fig:r2-s1-aks-json}
\end{figure}

\paragraph{Exécution du déploiement.}
Le cluster est créé depuis vm-gestion par la commande \texttt{aks-engine deploy}. La
figure~\ref{fig:r2-s1-aks-deploy} illustre l'exécution réussie de cette commande.

\begin{figure}[htbp]
    \centering
    \includegraphics[width=0.85\textwidth]{s1-aks-engine-deploy.png}
    \caption{Déploiement du cluster avec AKS Engine}
    \label{fig:r2-s1-aks-deploy}
\end{figure}

\paragraph{Vérification.}
La figure~\ref{fig:r2-s1-kubectl-nodes} présente la sortie de \texttt{kubectl get nodes} :
les trois nœuds du cluster sont en état \texttt{Ready}.

\begin{figure}[htbp]
    \centering
    \includegraphics[width=0.85\textwidth]{s1-kubectl-nodes.png}
    \caption{Vérification des nœuds Kubernetes (\texttt{kubectl get nodes})}
    \label{fig:r2-s1-kubectl-nodes}
\end{figure}


% ---------------------------------------------------------------------------
\subsection{Bilan du Sprint 1}
% ---------------------------------------------------------------------------

Le Sprint~1 fournit l'infrastructure réseau et compute nécessaire aux sprints suivants :
vm-gestion comme point d'ancrage DevOps, sous-réseau \texttt{aks} prêt pour les pods, et
cluster Kubernetes opérationnel. Le pipeline CI/CD (Sprint~2) et le déploiement applicatif
(Sprint~3) s'appuient sur ce socle.


% --- Sprint 2 : CI/CD ---------------------------------------------------------
% =============================================================================
% Release 2 — Sprint 2 : Intégration continue et déploiement continu (CI/CD)
% =============================================================================

\section[Sprint 2 : Intégration continue et déploiement continu (CI/CD)]%
{Sprint 2 : Intégration continue et déploiement continu (CI/CD)}

Ce sprint industrialise la validation du code, la construction des artefacts conteneurisés
et l'automatisation du provisionnement infrastructurel. Il s'appuie sur \textbf{Azure DevOps}
et un pipeline entièrement décrit en YAML (\texttt{azure-pipelines.yml}).


% ---------------------------------------------------------------------------
\subsection{Objectif du Sprint 2}
% ---------------------------------------------------------------------------

Le tableau~\ref{tab:r2-s2-objectifs} résume les objectifs du Sprint~2 par axe.

\begin{table}[H]
\centering
\caption{Objectifs du Sprint~2 par axe}
\label{tab:r2-s2-objectifs}
\footnotesize
\renewcommand{\arraystretch}{1.1}
\setlength{\tabcolsep}{6pt}
\begin{tabular}{|p{2.8cm}|p{9.2cm}|}
\hline
\textbf{Axe} & \textbf{Objectif} \\
\hline
Intégration continue &
  Pipeline déclenché sur \texttt{main} et \texttt{develop}, exécuté sur agent
  auto-hébergé (\texttt{ticketmanager-pool}) \\
\hline
Qualité &
  Analyse statique SonarQube ; tests automatisés backend (Jest) \\
\hline
Sécurité &
  Scan des images Docker (Trivy) ; analyse des dépendances (Snyk, optionnel) \\
\hline
Conteneurisation &
  Images multi-étapes \texttt{ticket-backend} et \texttt{ticket-frontend} \\
\hline
Registre &
  Publication des images vers Nexus Repository (\texttt{localhost:8082}) \\
\hline
IaC automatisé &
  Exécution Terraform (init, plan, apply) intégrée au pipeline \\
\hline
\end{tabular}
\end{table}

\paragraph{Résultat attendu.}
\begin{itemize}
  \item pipeline Azure DevOps vert de bout en bout (stages Build à Push minimum) ;
  \item images Docker versionnées dans Nexus (\texttt{docker-ticketmanager/}) ;
  \item rapport SonarQube pour le projet \texttt{ticket-manager} ;
  \item infrastructure Terraform synchronisée via le stage IaC.
\end{itemize}

\paragraph{Note sur les stages temporairement désactivés.}
Dans la version actuelle de \texttt{azure-pipelines.yml}, les stages \textbf{Test},
\textbf{Snyk} et \textbf{Trivy} sont commentés pour accélérer les itérations. Le squelette
du rapport prévoit leurs figures ; réactivez-les avant la soutenance si possible.


% ---------------------------------------------------------------------------
\subsection{Backlog du Sprint 2}
% ---------------------------------------------------------------------------

Le tableau~\ref{tab:r2-s2-backlog} présente le backlog planifié pour le Sprint~2.

\subsubsection{Epic — Chaîne d'intégration continue}

\begin{table}[H]
\centering
\caption{Backlog Sprint~2 — Epic intégration continue et déploiement continu}
\label{tab:r2-s2-backlog}
\footnotesize
\renewcommand{\arraystretch}{0.85}
\setlength{\tabcolsep}{4pt}
\begin{tabular}{|c|p{7.8cm}|c|c|}
\hline
\textbf{ID} & \textbf{User story} & \textbf{Priorité} & \textbf{Est. (j)} \\
\hline
R2-S2-C01 &
  En tant que développeur, je centralise le code source sur Azure Repos &
  H & 1 \\
\hline
R2-S2-C02 &
  En tant que développeur, j'exécute le pipeline sur un agent auto-hébergé &
  H & 2 \\
\hline
R2-S2-C03 &
  En tant que développeur, je déclenche automatiquement le pipeline à chaque push &
  H & 1 \\
\hline
R2-S2-C04 &
  En tant que développeur, je compile le backend NestJS et le frontend Next.js à chaque livraison &
  H & 2 \\
\hline
R2-S2-C05 &
  En tant que développeur, j'exécute les tests automatisés et publie la couverture &
  H & 2 \\
\hline
R2-S2-C06 &
  En tant que développeur, j'intègre l'analyse SonarQube dans le pipeline &
  H & 2 \\
\hline
R2-S2-C07 &
  En tant que développeur, je construis des images Docker prêtes pour la production &
  H & 2 \\
\hline
R2-S2-C08 &
  En tant que développeur, je scanne les images Docker avec Trivy avant publication &
  M & 1 \\
\hline
R2-S2-C09 &
  En tant que développeur, je publie les images validées dans Nexus Repository &
  H & 1 \\
\hline
R2-S2-C10 &
  En tant que développeur, j'automatise Terraform dans le pipeline avec validation manuelle &
  H & 2 \\
\hline
\end{tabular}
\end{table}


% ---------------------------------------------------------------------------
\subsection{Analyse — Conception du pipeline}
% ---------------------------------------------------------------------------

\subsubsection{Présentation d'Azure DevOps}

Azure DevOps regroupe les outils collaboratifs mobilisés dans TicketManager :
\begin{itemize}
  \item \textbf{Azure Boards} — planification agile et suivi des sprints ;
  \item \textbf{Azure Repos} — gestion de versions Git ;
  \item \textbf{Azure Pipelines} — exécution automatisée du pipeline CI/CD en YAML.
\end{itemize}

La figure~\ref{fig:r2-s2-ado-ui} présente l'interface principale d'Azure DevOps utilisée
dans le projet.

\begin{figure}[htbp]
    \centering
    \includegraphics[width=0.85\textwidth]{s2-azure-devops-ui.png}
    \caption{Interface principale d'Azure DevOps}
    \label{fig:r2-s2-ado-ui}
\end{figure}

\subsubsection{Structure déclarative du pipeline}

Le pipeline est organisé en \textbf{stages} séquentiels. La figure~\ref{fig:r2-s2-pipeline-stages}
présente la vue d'ensemble de ces étapes telles que définies dans \texttt{azure-pipelines.yml} :

\begin{enumerate}
  \item \textbf{Build} — compilation backend et frontend ;
  \item \textbf{CodeQuality} — analyse SonarQube ;
  \item \textbf{Docker} — construction des images ;
  \item \textbf{Push} — publication vers Nexus ;
  \item \textbf{Terraform} — provisionnement IaC ;
  \item \textbf{DeployK8s} — déploiement Kubernetes (Sprint~3, branche \texttt{main}) ;
  \item \textbf{DeployMonitoring} — Prometheus + Grafana via Ansible (Sprint~3).
\end{enumerate}

\begin{figure}[htbp]
    \centering
    \includegraphics[width=0.95\textwidth]{s2-pipeline-stages.png}
    \caption{Vue d'ensemble des stages du pipeline (\texttt{azure-pipelines.yml})}
    \label{fig:r2-s2-pipeline-stages}
\end{figure}

\subsubsection{Pool d'agents et connexion au dépôt}

Les tâches s'exécutent sur un \textbf{agent auto-hébergé} installé sur vm-gestion, regroupé
dans le pool \texttt{ticketmanager-pool}. La figure~\ref{fig:r2-s2-agent-pool} illustre la
configuration de ce pool d'agents. La figure~\ref{fig:r2-s2-repos} présente la structure du
dépôt Azure Repos, incluant les répertoires \texttt{backend}, \texttt{frontend},
\texttt{terraform} et \texttt{k8s}.

\begin{figure}[htbp]
    \centering
    \includegraphics[width=0.85\textwidth]{s2-agent-pool.png}
    \caption{Configuration du pool d'agents auto-hébergés}
    \label{fig:r2-s2-agent-pool}
\end{figure}

\begin{figure}[htbp]
    \centering
    \includegraphics[width=0.85\textwidth]{s2-azure-repos.png}
    \caption{Structure du dépôt Azure Repos (backend, frontend, terraform, k8s)}
    \label{fig:r2-s2-repos}
\end{figure}


% ---------------------------------------------------------------------------
\subsection{Réalisation}
% ---------------------------------------------------------------------------

\subsubsection{Compilation et tests}

\paragraph{Stage Build.}
Installation des dépendances npm (avec cache local sur vm-gestion), puis
\texttt{npm run build} pour le backend NestJS et le frontend Next.js.

\paragraph{Tests backend.}
Les tests unitaires et d'intégration s'exécutent avec Jest lorsque le stage Test est actif.
La figure~\ref{fig:r2-s2-tests-backend} présente le résultat de l'exécution des tests
automatisés du backend dans Azure DevOps.

\begin{figure}[htbp]
    \centering
    \includegraphics[width=0.85\textwidth]{s2-tests-backend.png}
    \caption{Exécution des tests automatisés du backend}
    \label{fig:r2-s2-tests-backend}
\end{figure}


\subsubsection{Analyse de la qualité du code (SonarQube)}

SonarQube analyse le code source backend (\texttt{backend/src}) et frontend
(\texttt{frontend}). La figure~\ref{fig:r2-s2-sonar-conn} présente la configuration de la
connexion entre SonarQube et Azure DevOps. La figure~\ref{fig:r2-s2-sonar-results} illustre
les résultats d'analyse du projet \texttt{ticket-manager}.

\begin{figure}[htbp]
    \centering
    \includegraphics[width=0.85\textwidth]{s2-sonar-connection.png}
    \caption{Connexion SonarQube — Azure DevOps}
    \label{fig:r2-s2-sonar-conn}
\end{figure}

\begin{figure}[htbp]
    \centering
    \includegraphics[width=0.85\textwidth]{s2-sonar-results.png}
    \caption{Résultats d'analyse SonarQube — projet TicketManager}
    \label{fig:r2-s2-sonar-results}
\end{figure}


\subsubsection{Construction des images Docker}

Deux images sont produites : \textbf{ticket-backend} (API NestJS, port 3001) et
\textbf{ticket-frontend} (Next.js, port 3000). La figure~\ref{fig:r2-s2-dockerfile-be}
présente le Dockerfile multi-étapes du backend. La figure~\ref{fig:r2-s2-dockerfile-fe}
détaille celui du frontend. La figure~\ref{fig:r2-s2-docker-build} montre la construction
réussie des images dans le stage Docker du pipeline.

\begin{figure}[htbp]
    \centering
    \includegraphics[width=0.85\textwidth]{s2-dockerfile-backend.png}
    \caption{Dockerfile du backend NestJS}
    \label{fig:r2-s2-dockerfile-be}
\end{figure}

\begin{figure}[htbp]
    \centering
    \includegraphics[width=0.85\textwidth]{s2-dockerfile-frontend.png}
    \caption{Dockerfile du frontend Next.js}
    \label{fig:r2-s2-dockerfile-fe}
\end{figure}

\begin{figure}[htbp]
    \centering
    \includegraphics[width=0.85\textwidth]{s2-docker-build.png}
    \caption{Construction des images dans le stage Docker}
    \label{fig:r2-s2-docker-build}
\end{figure}


\subsubsection{Scan de vulnérabilités (Trivy)}

Trivy analyse les images construites et signale les CVE par niveau de gravité. La
figure~\ref{fig:r2-s2-trivy-be} présente le rapport de scan pour l'image backend. La
figure~\ref{fig:r2-s2-trivy-fe} illustre celui de l'image frontend.

\begin{figure}[htbp]
    \centering
    \includegraphics[width=0.85\textwidth]{s2-trivy-backend.png}
    \caption{Scan Trivy — image backend}
    \label{fig:r2-s2-trivy-be}
\end{figure}

\begin{figure}[htbp]
    \centering
    \includegraphics[width=0.85\textwidth]{s2-trivy-frontend.png}
    \caption{Scan Trivy — image frontend}
    \label{fig:r2-s2-trivy-fe}
\end{figure}


\subsubsection{Publication vers Nexus Repository}

Nexus joue le rôle de \textbf{registre Docker privé} sur vm-gestion. La
figure~\ref{fig:r2-s2-nexus-repo} présente la création du dépôt Docker hosted. La
figure~\ref{fig:r2-s2-nexus-images} montre les images \texttt{ticket-backend} et
\texttt{ticket-frontend} stockées après exécution du pipeline.

\begin{figure}[htbp]
    \centering
    \includegraphics[width=0.85\textwidth]{s2-nexus-repo.png}
    \caption{Création du dépôt Docker hosted dans Nexus}
    \label{fig:r2-s2-nexus-repo}
\end{figure}

\begin{figure}[htbp]
    \centering
    \includegraphics[width=0.85\textwidth]{s2-nexus-images.png}
    \caption{Images \texttt{ticket-backend} et \texttt{ticket-frontend} dans Nexus}
    \label{fig:r2-s2-nexus-images}
\end{figure}


\subsubsection{Terraform automatisé dans le pipeline}

Le stage \textbf{Terraform IaC} exécute \texttt{terraform init}, \texttt{plan} et
\texttt{apply}. La figure~\ref{fig:r2-s2-tf-pipeline} présente l'exécution réussie de ce
stage dans Azure DevOps. La figure~\ref{fig:r2-s2-pipeline-run} illustre une exécution
complète du pipeline CI (stages Build à Push).

\begin{figure}[htbp]
    \centering
    \includegraphics[width=0.85\textwidth]{s2-terraform-pipeline.png}
    \caption{Exécution réussie du stage Terraform dans Azure DevOps}
    \label{fig:r2-s2-tf-pipeline}
\end{figure}

\begin{figure}[htbp]
    \centering
    \includegraphics[width=0.85\textwidth]{s2-pipeline-run-success.png}
    \caption{Exécution complète du pipeline CI (stages Build à Push)}
    \label{fig:r2-s2-pipeline-run}
\end{figure}


% ---------------------------------------------------------------------------
\subsection{Bilan du Sprint 2}
% ---------------------------------------------------------------------------

Le Sprint~2 met en place une chaîne d'intégration continue couvrant la compilation, la
qualité, la conteneurisation, la sécurité des images et la publication vers un registre
privé. Le déploiement effectif des pods applicatifs sur le cluster (stage
\textbf{DeployK8s}) et la supervision (stage \textbf{DeployMonitoring}) sont détaillés au
Sprint~3.


% --- Sprint 3 : Monitoring et CD ----------------------------------------------
% =============================================================================
% Release 2 — Sprint 3 : Monitoring et supervision de la plateforme
% =============================================================================

\section[Sprint 3 : Monitoring et supervision de la plateforme]%
{Sprint 3 : Monitoring et supervision de la plateforme}

Ce sprint finalise la Release~2 par le \textbf{déploiement continu} (CD) de l'application
sur Kubernetes et la mise en place de la \textbf{supervision} (Prometheus, Grafana) via
Ansible sur vm-gestion.


% ---------------------------------------------------------------------------
\subsection{Objectif du Sprint 3}
% ---------------------------------------------------------------------------

Le tableau~\ref{tab:r2-s3-objectifs} détaille les objectifs du Sprint~3 par axe.

\begin{table}[H]
\centering
\caption{Objectifs du Sprint~3 par axe}
\label{tab:r2-s3-objectifs}
\footnotesize
\renewcommand{\arraystretch}{1.1}
\setlength{\tabcolsep}{6pt}
\begin{tabular}{|p{2.8cm}|p{9.2cm}|}
\hline
\textbf{Axe} & \textbf{Objectif} \\
\hline
Déploiement continu &
  Déployer automatiquement backend, frontend et MongoDB in-cluster via le stage
  \texttt{DeployK8s} (branche \texttt{main}, approbation manuelle) \\
\hline
Manifests Kubernetes &
  Appliquer les fichiers YAML du répertoire \texttt{k8s/} (Deployments, Services,
  StatefulSet MongoDB, Secrets) \\
\hline
Registre privé &
  Configurer le pull d'images depuis Nexus (\texttt{10.1.1.4:8082}) sur les workers \\
\hline
Supervision &
  Déployer Prometheus et Grafana sur vm-gestion (\texttt{ansible/playbooks/monitoring.yml}) \\
\hline
Observabilité &
  Collecter les métriques des nœuds K8s, du backend NestJS et de MongoDB \\
\hline
Validation &
  Smoke tests post-déploiement (API, frontend, cibles Prometheus UP) \\
\hline
\end{tabular}
\end{table}

\paragraph{Résultat attendu.}
\begin{itemize}
  \item namespace \texttt{ticketmanager} avec pods \texttt{Running} ;
  \item application accessible via NodePort (frontend \texttt{30080}, backend \texttt{30081}) ;
  \item Prometheus accessible sur le port \texttt{9090}, Grafana sur \texttt{3030} ;
  \item dashboards opérationnels pour le cluster et l'application.
\end{itemize}

\paragraph{Choix architectural — MongoDB in-cluster.}
Contrairement aux projets précédents (chatbot) où MongoDB était exclusivement sur une VM
dédiée, TicketManager héberge la base de données de production dans un
\textbf{StatefulSet Kubernetes} avec volume persistant. La vm-db reste disponible pour
migration ou sauvegarde legacy (\texttt{ansible/playbooks/mongodb.yml}).


% ---------------------------------------------------------------------------
\subsection{Backlog du Sprint 3}
% ---------------------------------------------------------------------------

Le tableau~\ref{tab:r2-s3-backlog-deploy} couvre l'epic déploiement continu. Le
tableau~\ref{tab:r2-s3-backlog-monitoring} couvre l'epic supervision.

\subsubsection{Epic — Déploiement continu (Kubernetes)}

\begin{table}[H]
\centering
\caption{Backlog Sprint~3 — Epic déploiement continu Kubernetes}
\label{tab:r2-s3-backlog-deploy}
\footnotesize
\renewcommand{\arraystretch}{0.85}
\setlength{\tabcolsep}{4pt}
\begin{tabular}{|c|p{7.8cm}|c|c|}
\hline
\textbf{ID} & \textbf{User story} & \textbf{Priorité} & \textbf{Est. (j)} \\
\hline
R2-S3-D01 &
  En tant que développeur, je déploie le backend NestJS en pod Kubernetes &
  H & 2 \\
\hline
R2-S3-D02 &
  En tant que développeur, je déploie le frontend Next.js en pod Kubernetes &
  H & 2 \\
\hline
R2-S3-D03 &
  En tant que développeur, je déploie MongoDB in-cluster via un StatefulSet &
  H & 2 \\
\hline
R2-S3-D04 &
  En tant que développeur, je configure les secrets Kubernetes (JWT, MongoDB, Cloudinary) &
  H & 1 \\
\hline
R2-S3-D05 &
  En tant que développeur, j'automatise le déploiement via le stage CD du pipeline &
  H & 2 \\
\hline
R2-S3-D06 &
  En tant que développeur, je valide le déploiement par des smoke tests automatisés &
  H & 1 \\
\hline
\end{tabular}
\end{table}

\subsubsection{Epic — Monitoring et supervision}

\begin{table}[H]
\centering
\caption{Backlog Sprint~3 — Epic monitoring et supervision}
\label{tab:r2-s3-backlog-monitoring}
\footnotesize
\renewcommand{\arraystretch}{0.85}
\setlength{\tabcolsep}{4pt}
\begin{tabular}{|c|p{7.8cm}|c|c|}
\hline
\textbf{ID} & \textbf{User story} & \textbf{Priorité} & \textbf{Est. (j)} \\
\hline
R2-S3-M01 &
  En tant que développeur, je déploie Prometheus sur vm-gestion via Ansible &
  H & 2 \\
\hline
R2-S3-M02 &
  En tant que développeur, je déploie Grafana et je connecte Prometheus comme source de données &
  H & 2 \\
\hline
R2-S3-M03 &
  En tant que développeur, j'intègre le déploiement monitoring au pipeline Azure DevOps &
  M & 1 \\
\hline
R2-S3-M04 &
  En tant que développeur, je vérifie que les cibles Prometheus sont en état \texttt{UP} &
  H & 1 \\
\hline
\end{tabular}
\end{table}


% ---------------------------------------------------------------------------
\subsection{Analyse — Conception du déploiement Kubernetes}
% ---------------------------------------------------------------------------

\subsubsection{Architecture de déploiement}

Les composants applicatifs sont encapsulés dans des pods du namespace
\texttt{ticketmanager} :
\begin{itemize}
  \item \textbf{ticket-backend} — image Nexus, port 3001, secrets JWT et Cloudinary ;
  \item \textbf{ticket-frontend} — image Nexus, port 3000, URL API figée au build Docker ;
  \item \textbf{mongodb} — StatefulSet avec volume persistant pour les données tickets.
\end{itemize}

Les requêtes externes atteignent les services via \textbf{NodePort} sur les workers
(\texttt{30080} frontend, \texttt{30081} backend). La figure~\ref{fig:r2-s3-k8s-arch}
détaille l'architecture de déploiement Kubernetes (ressources du namespace
\texttt{ticketmanager}), en complément de la vue d'ensemble physique (figure~2.5).

\begin{figure}[htbp]
    \centering
    \includegraphics[width=0.92\textwidth]{architecture-deploiement-k8s.png}
    \caption{Architecture de déploiement Kubernetes — namespace \texttt{ticketmanager}}
    \label{fig:r2-s3-k8s-arch}
\end{figure}

\subsubsection{Chaîne de déploiement continu}

Le stage \texttt{DeployK8s} du pipeline :
\begin{enumerate}
  \item attend une \textbf{approbation manuelle} (\texttt{ManualValidation}) ;
  \item configure le registre insecure sur les workers ;
  \item exécute \texttt{k8s/deploy-pipeline.sh} (secrets, manifests, rollout, seed, smoke test).
\end{enumerate}


% ---------------------------------------------------------------------------
\subsection{Réalisation — Déploiement Kubernetes}
% ---------------------------------------------------------------------------

\subsubsection{Manifests et secrets}

La figure~\ref{fig:r2-s3-backend-yaml} présente le manifest de déploiement du backend
(\texttt{backend-deployment.yaml}). La figure~\ref{fig:r2-s3-frontend-yaml} illustre celui
du frontend (\texttt{frontend-deployment.yaml}).

\begin{figure}[htbp]
    \centering
    \includegraphics[width=0.85\textwidth]{s3-backend-deployment.png}
    \caption{Manifest de déploiement du backend (\texttt{backend-deployment.yaml})}
    \label{fig:r2-s3-backend-yaml}
\end{figure}

\begin{figure}[htbp]
    \centering
    \includegraphics[width=0.85\textwidth]{s3-frontend-deployment.png}
    \caption{Manifest de déploiement du frontend (\texttt{frontend-deployment.yaml})}
    \label{fig:r2-s3-frontend-yaml}
\end{figure}

\subsubsection{Exécution du déploiement}

La figure~\ref{fig:r2-s3-k8s-pipeline} présente le stage \textbf{DeployK8s} exécuté dans
Azure DevOps, incluant l'approbation manuelle et l'application des manifests.

\begin{figure}[htbp]
    \centering
    \includegraphics[width=0.85\textwidth]{s3-k8s-pipeline-stage.png}
    \caption{Stage DeployK8s dans Azure DevOps}
    \label{fig:r2-s3-k8s-pipeline}
\end{figure}

\subsubsection{Vérification des pods}

La figure~\ref{fig:r2-s3-pods} présente la sortie de \texttt{kubectl get pods -n ticketmanager}
: l'ensemble des pods applicatifs est en état \texttt{Running}.

\begin{figure}[htbp]
    \centering
    \includegraphics[width=0.85\textwidth]{s3-kubectl-pods.png}
    \caption{État des pods (\texttt{kubectl get pods -n ticketmanager})}
    \label{fig:r2-s3-pods}
\end{figure}

\subsubsection{Application accessible}

La figure~\ref{fig:r2-s3-app} illustre l'interface TicketManager accessible depuis un
navigateur après déploiement réussi sur le cluster.

\begin{figure}[htbp]
    \centering
    \includegraphics[width=0.85\textwidth]{s3-app-browser.png}
    \caption{Interface TicketManager accessible après déploiement}
    \label{fig:r2-s3-app}
\end{figure}


% ---------------------------------------------------------------------------
\subsection{Réalisation — Monitoring et supervision}
% ---------------------------------------------------------------------------

\subsubsection{Déploiement Ansible (Prometheus + Grafana)}

Le playbook \texttt{ansible/playbooks/monitoring.yml} déploie Prometheus (port 9090) et
Grafana (port 3030) sur vm-gestion. La figure~\ref{fig:r2-s3-ansible} présente l'exécution
réussie de ce playbook.

\begin{figure}[htbp]
    \centering
    \includegraphics[width=0.85\textwidth]{s3-ansible-monitoring.png}
    \caption{Exécution du playbook \texttt{monitoring.yml}}
    \label{fig:r2-s3-ansible}
\end{figure}

\subsubsection{Configuration Prometheus}

Le fichier \texttt{ansible/templates/prometheus.yml.j2} définit les cibles surveillées.
La figure~\ref{fig:r2-s3-prom-targets} présente l'interface Prometheus (\texttt{/targets})
avec l'ensemble des cibles en état \texttt{UP}.

\begin{figure}[htbp]
    \centering
    \includegraphics[width=0.85\textwidth]{s3-prometheus-targets.png}
    \caption{Cibles actives dans l'interface Prometheus (\texttt{/targets})}
    \label{fig:r2-s3-prom-targets}
\end{figure}

\subsubsection{Dashboards Grafana}

La figure~\ref{fig:r2-s3-grafana-ds} présente la configuration de Prometheus comme source
de données dans Grafana. La figure~\ref{fig:r2-s3-grafana-dash} illustre un dashboard de
supervision regroupant les métriques du cluster et de l'application backend.

\begin{figure}[htbp]
    \centering
    \includegraphics[width=0.85\textwidth]{s3-grafana-datasource.png}
    \caption{Configuration de Prometheus comme source Grafana}
    \label{fig:r2-s3-grafana-ds}
\end{figure}

\begin{figure}[htbp]
    \centering
    \includegraphics[width=0.85\textwidth]{s3-grafana-dashboard.png}
    \caption{Dashboard de supervision — cluster et application}
    \label{fig:r2-s3-grafana-dash}
\end{figure}

\subsubsection{Intégration pipeline}

La figure~\ref{fig:r2-s3-mon-pipeline} présente le stage \textbf{DeployMonitoring} du
pipeline Azure DevOps, exécutant automatiquement le playbook Ansible après un déploiement
K8s réussi.

\begin{figure}[htbp]
    \centering
    \includegraphics[width=0.85\textwidth]{s3-monitoring-pipeline.png}
    \caption{Stage DeployMonitoring dans Azure DevOps}
    \label{fig:r2-s3-mon-pipeline}
\end{figure}


% ---------------------------------------------------------------------------
\subsection{Tests et validation}
% ---------------------------------------------------------------------------

\subsubsection{Smoke tests post-déploiement}

Le script \texttt{k8s/post-deploy.sh} vérifie la disponibilité de l'API, du frontend et
des outils de supervision. Le tableau~\ref{tab:r2-s3-validation} résume les critères de
validation retenus pour le Sprint~3.

\begin{table}[H]
\centering
\caption{Critères de validation du Sprint~3}
\label{tab:r2-s3-validation}
\footnotesize
\renewcommand{\arraystretch}{1.1}
\setlength{\tabcolsep}{6pt}
\begin{tabular}{|l|l|l|}
\hline
\textbf{Composant} & \textbf{URL / commande} & \textbf{Critère de succès} \\
\hline
Backend & \texttt{curl http://<worker>:30081/} & HTTP 200 \\
\hline
Frontend & \texttt{http://<worker>:30080} & Page de connexion affichée \\
\hline
Pods & \texttt{kubectl get pods -n ticketmanager} & Tous \texttt{Running} \\
\hline
Prometheus & \texttt{curl http://10.1.1.4:9090/-/ready} & \texttt{Ready} \\
\hline
Grafana & \texttt{http://10.1.1.4:3030/login} & HTTP 200 \\
\hline
\end{tabular}
\end{table}


% ---------------------------------------------------------------------------
\subsection{Bilan du Sprint 3}
% ---------------------------------------------------------------------------

Le Sprint~3 clôt la Release~2 en rendant TicketManager déployable et observable en
environnement cloud. L'application tourne dans Kubernetes ; la supervision Prometheus/Grafana
permet de détecter les anomalies de charge, de disponibilité et de performance avant
intervention manuelle.


% --- Bilan Release 2 ----------------------------------------------------------
\subsection{Bilan de la Release~2}

À l'issue des trois sprints, la plateforme \textit{TicketManager} dispose d'une
infrastructure documentée et reproductible sur Azure Stack Hub, d'une chaîne CI/CD
automatisée (compilation, qualité, conteneurisation, publication Nexus, provisionnement
IaC) et d'un déploiement applicatif supervisé (pods Kubernetes, Prometheus, Grafana).

Les indicateurs de réussite retenus pour cette release sont :
\begin{itemize}
  \item exécution réussie du pipeline Azure DevOps sur la branche \texttt{main} ;
  \item pods du namespace \texttt{ticketmanager} en état \texttt{Running} ;
  \item accessibilité de l'interface web et de l'API depuis le réseau autorisé ;
  \item cibles Prometheus en état \texttt{UP} et dashboards Grafana opérationnels.
\end{itemize}

Les évolutions ultérieures (hors Release~2) peuvent porter sur le renforcement des tests
automatisés dans le pipeline, la réactivation des stages Snyk/Trivy, l'exposition via
Ingress Controller et la haute disponibilité multi-réplicas.

%
% Convention : les \caption des tableaux sont placés AU-DESSUS du tabular.
% Les figures de réalisation sont citées dans le texte via \ref{fig:...}.
%
% =============================================================================

\section*{Release 2 : Industrialisation et déploiement de la plateforme}
\addcontentsline{toc}{section}{Release 2 : Industrialisation et déploiement de la plateforme}

La \textbf{Release~2} prolonge le socle applicatif livré en Release~1 (conception et
développement de \textit{TicketManager}). Elle vise à industrialiser le cycle de vie de la
plateforme sur \textbf{Azure Stack Hub} : provisionnement reproductible de l'infrastructure,
chaîne d'intégration et de déploiement continus (CI/CD), déploiement conteneurisé sur
Kubernetes et supervision de l'environnement de production.

\paragraph{Organisation en trois sprints.}
Chaque sprint suit la même trame pédagogique que la Release~1 :
\begin{enumerate}
  \item \textbf{objectifs} fixés et périmètre ;
  \item \textbf{backlog} planifié (user stories) ;
  \item \textbf{analyse} (conception technique, architecture) ;
  \item \textbf{réalisation} (mise en œuvre, captures d'écran, validation).
\end{enumerate}

\paragraph{Codification des priorités (backlog).}
Dans les tableaux de backlog, la colonne \textbf{Priorité} utilise les codes suivants :
\textbf{H} (haute), \textbf{M} (moyenne) et \textbf{B} (basse).

Le tableau~\ref{tab:r2-overview} synthétise les trois sprints de la release.

\begin{table}[H]
\centering
\caption{Vue d'ensemble de la Release~2}
\label{tab:r2-overview}
\footnotesize
\renewcommand{\arraystretch}{1.15}
\setlength{\tabcolsep}{6pt}
\begin{tabular}{|p{1.2cm}|p{4.2cm}|p{7.8cm}|}
\hline
\textbf{Sprint} & \textbf{Intitulé} & \textbf{Livrables principaux} \\
\hline
Sprint~1 & Mise en place de l'infrastructure &
  Terraform (VNet, NSG, VMs), AKS Engine (cluster K8s), plan d'adressage \\
\hline
Sprint~2 & Intégration continue et déploiement continu (CI/CD) &
  Pipeline Azure DevOps, SonarQube, Docker, Nexus, Trivy, Terraform automatisé \\
\hline
Sprint~3 & Monitoring et supervision de la plateforme &
  Déploiement K8s (CD), MongoDB in-cluster, Prometheus, Grafana (Ansible) \\
\hline
\end{tabular}
\end{table}

\paragraph{Périmètre par rapport à la Release~1.}
La Release~1 couvre le développement fonctionnel (NestJS, Next.js, MongoDB locale,
Cloudinary, notifications WebSocket). La Release~2 ne modifie pas le métier applicatif :
elle prépare et automatise l'environnement cloud, la livraison des images Docker et le
déploiement des composants \texttt{ticket-backend}, \texttt{ticket-frontend} et
\texttt{mongodb} dans le namespace Kubernetes \texttt{ticketmanager}.

\paragraph{Fichiers sources du dépôt.}
\begin{itemize}
  \item Infrastructure : répertoire \texttt{terraform/} ;
  \item Pipeline : \texttt{azure-pipelines.yml} ;
  \item Déploiement Kubernetes : répertoire \texttt{k8s/} ;
  \item Automatisation : répertoire \texttt{ansible/} ;
  \item Documentation Release~1 (socle applicatif) : \texttt{docs/release-1/}.
\end{itemize}

% --- Sprint 1 : Infrastructure ------------------------------------------------
% =============================================================================
% Release 2 — Sprint 1 : Mise en place de l'infrastructure
% =============================================================================

\section[Sprint 1 : Mise en place de l'infrastructure]%
{Sprint 1 : Mise en place de l'infrastructure}

Ce sprint pose les fondations cloud de \textit{TicketManager} sur Azure Stack Hub.
L'objectif est de disposer d'un réseau segmenté, de machines virtuelles d'administration
et de base de données, ainsi que d'un cluster Kubernetes opérationnel, le tout
provisionné de manière reproductible.


% ---------------------------------------------------------------------------
\subsection{Objectif du Sprint 1}
% ---------------------------------------------------------------------------

Le Sprint~1 succède au socle applicatif de la Release~1. Il se concentre sur
\textbf{l'approvisionnement de l'infrastructure} et la \textbf{préparation du cluster}
destiné à héberger l'application conteneurisée. Le tableau~\ref{tab:r2-s1-objectifs} détaille
les objectifs par axe.

\begin{table}[H]
\centering
\caption{Objectifs du Sprint~1 par axe}
\label{tab:r2-s1-objectifs}
\footnotesize
\renewcommand{\arraystretch}{1.1}
\setlength{\tabcolsep}{6pt}
\begin{tabular}{|p{2.8cm}|p{9.2cm}|}
\hline
\textbf{Axe} & \textbf{Objectif} \\
\hline
Infrastructure as Code &
  Décrire et provisionner le réseau, la sécurité et les VMs par Terraform
  (fournisseur \textit{azurestack}) \\
\hline
Segmentation réseau &
  Structurer le VNet \texttt{10.1.0.0/16} en sous-réseaux \texttt{mgmt}, \texttt{db}
  et \texttt{aks}, chacun protégé par un NSG \\
\hline
Machines virtuelles &
  Déployer \textbf{vm-gestion} (administration, agent DevOps) et \textbf{vm-db}
  (MongoDB legacy / migration) \\
\hline
Orchestration &
  Créer le cluster Kubernetes (1 master + 2 workers) via \textbf{AKS Engine} \\
\hline
Vérification &
  Contrôler les ressources dans le portail Azure Stack Hub et via \texttt{kubectl} \\
\hline
\end{tabular}
\end{table}

\paragraph{Résultat attendu.}
À l'issue du sprint :
\begin{itemize}
  \item le réseau virtuel \texttt{ticketmanager-iac-vnet} et les trois sous-réseaux sont
        créés ;
  \item \textbf{vm-gestion} (\texttt{10.1.1.4}) est accessible en SSH et héberge l'agent
        Azure DevOps, Docker, Ansible et AKS Engine ;
  \item \textbf{vm-db} (\texttt{10.1.2.4}) est joignable en IP privée depuis \texttt{mgmt}
        et \texttt{aks} ;
  \item le cluster Kubernetes affiche trois nœuds \texttt{Ready}
        (\texttt{kubectl get nodes}).
\end{itemize}

\paragraph{Périmètre.}
Le tableau~\ref{tab:r2-s1-perimetre} précise le contenu du sprint par rapport aux livraisons
ultérieures.

\begin{table}[H]
\centering
\caption{Périmètre du Sprint~1}
\label{tab:r2-s1-perimetre}
\footnotesize
\renewcommand{\arraystretch}{1.1}
\setlength{\tabcolsep}{6pt}
\begin{tabular}{|p{5.5cm}|c|c|}
\hline
\textbf{Élément} & \textbf{Sprint~1} & \textbf{Sprint ultérieur} \\
\hline
Socle applicatif (Release~1) & Livré & — \\
\hline
Terraform — VNet, NSG, NIC, VMs & Oui & — \\
\hline
AKS Engine — cluster K8s & Oui & — \\
\hline
Pipeline Azure DevOps & — & Sprint~2 \\
\hline
Déploiement pods applicatifs & — & Sprint~3 \\
\hline
Prometheus / Grafana & — & Sprint~3 \\
\hline
\end{tabular}
\end{table}


% ---------------------------------------------------------------------------
\subsection{Backlog du Sprint 1}
% ---------------------------------------------------------------------------

Le backlog du Sprint~1 est organisé en deux epics : infrastructure Terraform et cluster
Kubernetes. Les tableaux~\ref{tab:r2-s1-backlog-infra} et~\ref{tab:r2-s1-backlog-k8s} listent
les user stories planifiées.

\subsubsection{Epic — Infrastructure cloud (Terraform)}

\begin{table}[H]
\centering
\caption{Backlog Sprint~1 — Epic infrastructure Terraform}
\label{tab:r2-s1-backlog-infra}
\footnotesize
\renewcommand{\arraystretch}{0.9}
\setlength{\tabcolsep}{4pt}
\begin{tabular}{|c|p{7.8cm}|c|c|}
\hline
\textbf{ID} & \textbf{User story} & \textbf{Priorité} & \textbf{Est. (j)} \\
\hline
R2-S1-I01 &
  En tant que développeur, je provisionne l'infrastructure par Terraform sur Azure Stack Hub &
  H & 3 \\
\hline
R2-S1-I02 &
  En tant que développeur, je segmente le réseau en trois sous-réseaux dédiés (mgmt, db, aks) &
  H & 2 \\
\hline
R2-S1-I03 &
  En tant que développeur, je configure les groupes de sécurité réseau par sous-réseau &
  H & 2 \\
\hline
R2-S1-I04 &
  En tant que développeur, je dispose d'une VM d'administration (vm-gestion) pour orchestrer les déploiements &
  H & 2 \\
\hline
R2-S1-I05 &
  En tant que développeur, je déploie une VM base de données (vm-db) pour MongoDB Legacy &
  M & 1 \\
\hline
R2-S1-I06 &
  En tant que développeur, j'associe chaque VM à une interface réseau sécurisée &
  H & 1 \\
\hline
\end{tabular}
\end{table}

\subsubsection{Epic — Cluster Kubernetes (AKS Engine)}

\begin{table}[H]
\centering
\caption{Backlog Sprint~1 — Epic cluster Kubernetes}
\label{tab:r2-s1-backlog-k8s}
\footnotesize
\renewcommand{\arraystretch}{0.9}
\setlength{\tabcolsep}{4pt}
\begin{tabular}{|c|p{7.8cm}|c|c|}
\hline
\textbf{ID} & \textbf{User story} & \textbf{Priorité} & \textbf{Est. (j)} \\
\hline
R2-S1-K01 &
  En tant que développeur, je prépare le fichier de configuration AKS Engine (\texttt{kubernetes.json}) &
  H & 2 \\
\hline
R2-S1-K02 &
  En tant que développeur, je déploie le cluster Kubernetes (1 master + 2 workers) &
  H & 3 \\
\hline
R2-S1-K03 &
  En tant que développeur, je vérifie l'état des nœuds du cluster (\texttt{Ready}) &
  H & 1 \\
\hline
R2-S1-K04 &
  En tant que développeur, je configure l'accès \texttt{kubeconfig} depuis vm-gestion &
  H & 1 \\
\hline
\end{tabular}
\end{table}


% ---------------------------------------------------------------------------
\subsection{Conception de l'infrastructure}
% ---------------------------------------------------------------------------

\subsubsection{Plan d'adressage réseau}

Le tableau~\ref{tab:r2-s1-adressage} résume le plan d'adressage retenu pour l'infrastructure
TicketManager. La figure~\ref{fig:r2-s1-plan-adressage} en propose une représentation
schématique (sous-réseaux, adresses IP et séparation Terraform / AKS Engine).

\begin{table}[H]
\centering
\caption{Plan d'adressage réseau TicketManager}
\label{tab:r2-s1-adressage}
\footnotesize
\renewcommand{\arraystretch}{1.1}
\setlength{\tabcolsep}{6pt}
\begin{tabular}{|l|l|}
\hline
\textbf{Élément} & \textbf{Valeur} \\
\hline
Réseau virtuel (VNet) & \texttt{10.1.0.0/16} \\
\hline
Subnet administration (\texttt{mgmt}) & \texttt{10.1.1.0/24} \\
\hline
Subnet base de données (\texttt{db}) & \texttt{10.1.2.0/24} \\
\hline
Subnet Kubernetes (\texttt{aks}) & \texttt{10.1.3.0/24} \\
\hline
IP vm-gestion (privée) & \texttt{10.1.1.4} \\
\hline
IP vm-db (privée) & \texttt{10.1.2.4} \\
\hline
Workers K8s (exemple) & \texttt{10.1.3.4}, \texttt{10.1.3.5} \\
\hline
\end{tabular}
\end{table}

\begin{figure}[htbp]
    \centering
    \includegraphics[width=0.95\textwidth]{plan-adressage-reseau-sprint1.png}
    \caption{Plan d'adressage réseau TicketManager — vue schématique (draw.io)}
    \label{fig:r2-s1-plan-adressage}
\end{figure}

\subsubsection{Architecture globale}

L'architecture repose sur une organisation \textbf{hybride} : couche réseau et compute
provisionnée par Terraform, complétée par un cluster Kubernetes déployé via AKS Engine
sur le sous-réseau \texttt{aks}. La figure~\ref{fig:r2-s1-arch-infra} en donne une vue
d'ensemble.

\begin{figure}[htbp]
    \centering
    \includegraphics[width=0.95\textwidth]{architecture-infrastructure-sprint1.png}
    \caption{Architecture globale de l'infrastructure TicketManager (Sprint~1)}
    \label{fig:r2-s1-arch-infra}
\end{figure}

\paragraph{Composants provisionnés par Terraform.}
\begin{itemize}
  \item \textbf{Réseau virtuel et sous-réseaux} : connectivité interne entre VMs et cluster ;
  \item \textbf{Groupes de sécurité réseau (NSG)} : filtrage SSH, MongoDB, trafic interne ;
  \item \textbf{Interfaces réseau (NIC)} : rattachement des VMs au VNet ;
  \item \textbf{Machines virtuelles} : vm-gestion et vm-db
        (fichiers \texttt{vm-management.tf}, \texttt{vm-db.tf}).
\end{itemize}

\paragraph{Cluster Kubernetes (hors Terraform).}
Terraform prépare le sous-réseau \texttt{aks} et les règles NSG associées ; le déploiement
effectif du plan de contrôle et des workers est réalisé par \textbf{AKS Engine}, outil
open-source adapté à Azure Stack Hub.


% ---------------------------------------------------------------------------
\subsection{Réalisation}
% ---------------------------------------------------------------------------

\subsubsection{Infrastructure as Code avec Terraform}

Terraform adopte une approche \textbf{déclarative} : l'état cible des ressources est décrit
dans des fichiers de configuration ; le moteur calcule et applique les écarts. Le cycle
repose sur quatre commandes : \texttt{terraform init}, \texttt{validate}, \texttt{plan} et
\texttt{apply}.

La figure~\ref{fig:r2-s1-tf-init} présente le résultat de l'initialisation de l'environnement
Terraform. La figure~\ref{fig:r2-s1-tf-validate} confirme la validité syntaxique de la
configuration. La figure~\ref{fig:r2-s1-tf-plan} illustre le plan d'exécution généré avant
application. Enfin, la figure~\ref{fig:r2-s1-tf-apply} montre les sorties du déploiement,
notamment les adresses IP attribuées aux machines virtuelles.

\begin{figure}[htbp]
    \centering
    \includegraphics[width=0.85\textwidth]{s1-terraform-init.png}
    \caption{Initialisation Terraform (\texttt{terraform init})}
    \label{fig:r2-s1-tf-init}
\end{figure}

\begin{figure}[htbp]
    \centering
    \includegraphics[width=0.85\textwidth]{s1-terraform-validate.png}
    \caption{Validation Terraform (\texttt{terraform validate})}
    \label{fig:r2-s1-tf-validate}
\end{figure}

\begin{figure}[htbp]
    \centering
    \includegraphics[width=0.85\textwidth]{s1-terraform-plan.png}
    \caption{Plan d'exécution Terraform (\texttt{terraform plan})}
    \label{fig:r2-s1-tf-plan}
\end{figure}

\begin{figure}[htbp]
    \centering
    \includegraphics[width=0.85\textwidth]{s1-terraform-apply.png}
    \caption{Application Terraform — sorties (IPs vm-gestion, vm-db)}
    \label{fig:r2-s1-tf-apply}
\end{figure}

\paragraph{Vérification portail Azure Stack Hub.}
Les figures~\ref{fig:r2-s1-vnet},~\ref{fig:r2-s1-nsg} et~\ref{fig:r2-s1-vms} confirment la
création des ressources dans le groupe \texttt{AmalBH-RG}. La figure~\ref{fig:r2-s1-vnet}
présente le réseau virtuel et les sous-réseaux. La figure~\ref{fig:r2-s1-nsg} détaille les
groupes de sécurité réseau. La figure~\ref{fig:r2-s1-vms} liste les machines virtuelles
déployées.

\begin{figure}[htbp]
    \centering
    \includegraphics[width=0.85\textwidth]{s1-vnet-subnets.png}
    \caption{Réseau virtuel et sous-réseaux dans le portail}
    \label{fig:r2-s1-vnet}
\end{figure}

\begin{figure}[htbp]
    \centering
    \includegraphics[width=0.85\textwidth]{s1-nsg.png}
    \caption{Groupes de sécurité réseau}
    \label{fig:r2-s1-nsg}
\end{figure}

\begin{figure}[htbp]
    \centering
    \includegraphics[width=0.85\textwidth]{s1-vms-portal.png}
    \caption{Machines virtuelles déployées}
    \label{fig:r2-s1-vms}
\end{figure}


\subsubsection{Déploiement du cluster Kubernetes avec AKS Engine}

\paragraph{Configuration du cluster.}
Le fichier \texttt{kubernetes.json} (modèle API AKS Engine) spécifie la version de
Kubernetes, le plugin réseau (\texttt{kubenet}), le runtime (\texttt{containerd}), la
topologie des nœuds et l'authentification SSH. La figure~\ref{fig:r2-s1-aks-json} présente
un extrait de ce fichier de configuration.

\begin{figure}[htbp]
    \centering
    \includegraphics[width=0.85\textwidth]{s1-aks-engine-json.png}
    \caption{Extrait du fichier de configuration AKS Engine}
    \label{fig:r2-s1-aks-json}
\end{figure}

\paragraph{Exécution du déploiement.}
Le cluster est créé depuis vm-gestion par la commande \texttt{aks-engine deploy}. La
figure~\ref{fig:r2-s1-aks-deploy} illustre l'exécution réussie de cette commande.

\begin{figure}[htbp]
    \centering
    \includegraphics[width=0.85\textwidth]{s1-aks-engine-deploy.png}
    \caption{Déploiement du cluster avec AKS Engine}
    \label{fig:r2-s1-aks-deploy}
\end{figure}

\paragraph{Vérification.}
La figure~\ref{fig:r2-s1-kubectl-nodes} présente la sortie de \texttt{kubectl get nodes} :
les trois nœuds du cluster sont en état \texttt{Ready}.

\begin{figure}[htbp]
    \centering
    \includegraphics[width=0.85\textwidth]{s1-kubectl-nodes.png}
    \caption{Vérification des nœuds Kubernetes (\texttt{kubectl get nodes})}
    \label{fig:r2-s1-kubectl-nodes}
\end{figure}


% ---------------------------------------------------------------------------
\subsection{Bilan du Sprint 1}
% ---------------------------------------------------------------------------

Le Sprint~1 fournit l'infrastructure réseau et compute nécessaire aux sprints suivants :
vm-gestion comme point d'ancrage DevOps, sous-réseau \texttt{aks} prêt pour les pods, et
cluster Kubernetes opérationnel. Le pipeline CI/CD (Sprint~2) et le déploiement applicatif
(Sprint~3) s'appuient sur ce socle.


% --- Sprint 2 : CI/CD ---------------------------------------------------------
% =============================================================================
% Release 2 — Sprint 2 : Intégration continue et déploiement continu (CI/CD)
% =============================================================================

\section[Sprint 2 : Intégration continue et déploiement continu (CI/CD)]%
{Sprint 2 : Intégration continue et déploiement continu (CI/CD)}

Ce sprint industrialise la validation du code, la construction des artefacts conteneurisés
et l'automatisation du provisionnement infrastructurel. Il s'appuie sur \textbf{Azure DevOps}
et un pipeline entièrement décrit en YAML (\texttt{azure-pipelines.yml}).


% ---------------------------------------------------------------------------
\subsection{Objectif du Sprint 2}
% ---------------------------------------------------------------------------

Le tableau~\ref{tab:r2-s2-objectifs} résume les objectifs du Sprint~2 par axe.

\begin{table}[H]
\centering
\caption{Objectifs du Sprint~2 par axe}
\label{tab:r2-s2-objectifs}
\footnotesize
\renewcommand{\arraystretch}{1.1}
\setlength{\tabcolsep}{6pt}
\begin{tabular}{|p{2.8cm}|p{9.2cm}|}
\hline
\textbf{Axe} & \textbf{Objectif} \\
\hline
Intégration continue &
  Pipeline déclenché sur \texttt{main} et \texttt{develop}, exécuté sur agent
  auto-hébergé (\texttt{ticketmanager-pool}) \\
\hline
Qualité &
  Analyse statique SonarQube ; tests automatisés backend (Jest) \\
\hline
Sécurité &
  Scan des images Docker (Trivy) ; analyse des dépendances (Snyk, optionnel) \\
\hline
Conteneurisation &
  Images multi-étapes \texttt{ticket-backend} et \texttt{ticket-frontend} \\
\hline
Registre &
  Publication des images vers Nexus Repository (\texttt{localhost:8082}) \\
\hline
IaC automatisé &
  Exécution Terraform (init, plan, apply) intégrée au pipeline \\
\hline
\end{tabular}
\end{table}

\paragraph{Résultat attendu.}
\begin{itemize}
  \item pipeline Azure DevOps vert de bout en bout (stages Build à Push minimum) ;
  \item images Docker versionnées dans Nexus (\texttt{docker-ticketmanager/}) ;
  \item rapport SonarQube pour le projet \texttt{ticket-manager} ;
  \item infrastructure Terraform synchronisée via le stage IaC.
\end{itemize}

\paragraph{Note sur les stages temporairement désactivés.}
Dans la version actuelle de \texttt{azure-pipelines.yml}, les stages \textbf{Test},
\textbf{Snyk} et \textbf{Trivy} sont commentés pour accélérer les itérations. Le squelette
du rapport prévoit leurs figures ; réactivez-les avant la soutenance si possible.


% ---------------------------------------------------------------------------
\subsection{Backlog du Sprint 2}
% ---------------------------------------------------------------------------

Le tableau~\ref{tab:r2-s2-backlog} présente le backlog planifié pour le Sprint~2.

\subsubsection{Epic — Chaîne d'intégration continue}

\begin{table}[H]
\centering
\caption{Backlog Sprint~2 — Epic intégration continue et déploiement continu}
\label{tab:r2-s2-backlog}
\footnotesize
\renewcommand{\arraystretch}{0.85}
\setlength{\tabcolsep}{4pt}
\begin{tabular}{|c|p{7.8cm}|c|c|}
\hline
\textbf{ID} & \textbf{User story} & \textbf{Priorité} & \textbf{Est. (j)} \\
\hline
R2-S2-C01 &
  En tant que développeur, je centralise le code source sur Azure Repos &
  H & 1 \\
\hline
R2-S2-C02 &
  En tant que développeur, j'exécute le pipeline sur un agent auto-hébergé &
  H & 2 \\
\hline
R2-S2-C03 &
  En tant que développeur, je déclenche automatiquement le pipeline à chaque push &
  H & 1 \\
\hline
R2-S2-C04 &
  En tant que développeur, je compile le backend NestJS et le frontend Next.js à chaque livraison &
  H & 2 \\
\hline
R2-S2-C05 &
  En tant que développeur, j'exécute les tests automatisés et publie la couverture &
  H & 2 \\
\hline
R2-S2-C06 &
  En tant que développeur, j'intègre l'analyse SonarQube dans le pipeline &
  H & 2 \\
\hline
R2-S2-C07 &
  En tant que développeur, je construis des images Docker prêtes pour la production &
  H & 2 \\
\hline
R2-S2-C08 &
  En tant que développeur, je scanne les images Docker avec Trivy avant publication &
  M & 1 \\
\hline
R2-S2-C09 &
  En tant que développeur, je publie les images validées dans Nexus Repository &
  H & 1 \\
\hline
R2-S2-C10 &
  En tant que développeur, j'automatise Terraform dans le pipeline avec validation manuelle &
  H & 2 \\
\hline
\end{tabular}
\end{table}


% ---------------------------------------------------------------------------
\subsection{Analyse — Conception du pipeline}
% ---------------------------------------------------------------------------

\subsubsection{Présentation d'Azure DevOps}

Azure DevOps regroupe les outils collaboratifs mobilisés dans TicketManager :
\begin{itemize}
  \item \textbf{Azure Boards} — planification agile et suivi des sprints ;
  \item \textbf{Azure Repos} — gestion de versions Git ;
  \item \textbf{Azure Pipelines} — exécution automatisée du pipeline CI/CD en YAML.
\end{itemize}

La figure~\ref{fig:r2-s2-ado-ui} présente l'interface principale d'Azure DevOps utilisée
dans le projet.

\begin{figure}[htbp]
    \centering
    \includegraphics[width=0.85\textwidth]{s2-azure-devops-ui.png}
    \caption{Interface principale d'Azure DevOps}
    \label{fig:r2-s2-ado-ui}
\end{figure}

\subsubsection{Structure déclarative du pipeline}

Le pipeline est organisé en \textbf{stages} séquentiels. La figure~\ref{fig:r2-s2-pipeline-stages}
présente la vue d'ensemble de ces étapes telles que définies dans \texttt{azure-pipelines.yml} :

\begin{enumerate}
  \item \textbf{Build} — compilation backend et frontend ;
  \item \textbf{CodeQuality} — analyse SonarQube ;
  \item \textbf{Docker} — construction des images ;
  \item \textbf{Push} — publication vers Nexus ;
  \item \textbf{Terraform} — provisionnement IaC ;
  \item \textbf{DeployK8s} — déploiement Kubernetes (Sprint~3, branche \texttt{main}) ;
  \item \textbf{DeployMonitoring} — Prometheus + Grafana via Ansible (Sprint~3).
\end{enumerate}

\begin{figure}[htbp]
    \centering
    \includegraphics[width=0.95\textwidth]{s2-pipeline-stages.png}
    \caption{Vue d'ensemble des stages du pipeline (\texttt{azure-pipelines.yml})}
    \label{fig:r2-s2-pipeline-stages}
\end{figure}

\subsubsection{Pool d'agents et connexion au dépôt}

Les tâches s'exécutent sur un \textbf{agent auto-hébergé} installé sur vm-gestion, regroupé
dans le pool \texttt{ticketmanager-pool}. La figure~\ref{fig:r2-s2-agent-pool} illustre la
configuration de ce pool d'agents. La figure~\ref{fig:r2-s2-repos} présente la structure du
dépôt Azure Repos, incluant les répertoires \texttt{backend}, \texttt{frontend},
\texttt{terraform} et \texttt{k8s}.

\begin{figure}[htbp]
    \centering
    \includegraphics[width=0.85\textwidth]{s2-agent-pool.png}
    \caption{Configuration du pool d'agents auto-hébergés}
    \label{fig:r2-s2-agent-pool}
\end{figure}

\begin{figure}[htbp]
    \centering
    \includegraphics[width=0.85\textwidth]{s2-azure-repos.png}
    \caption{Structure du dépôt Azure Repos (backend, frontend, terraform, k8s)}
    \label{fig:r2-s2-repos}
\end{figure}


% ---------------------------------------------------------------------------
\subsection{Réalisation}
% ---------------------------------------------------------------------------

\subsubsection{Compilation et tests}

\paragraph{Stage Build.}
Installation des dépendances npm (avec cache local sur vm-gestion), puis
\texttt{npm run build} pour le backend NestJS et le frontend Next.js.

\paragraph{Tests backend.}
Les tests unitaires et d'intégration s'exécutent avec Jest lorsque le stage Test est actif.
La figure~\ref{fig:r2-s2-tests-backend} présente le résultat de l'exécution des tests
automatisés du backend dans Azure DevOps.

\begin{figure}[htbp]
    \centering
    \includegraphics[width=0.85\textwidth]{s2-tests-backend.png}
    \caption{Exécution des tests automatisés du backend}
    \label{fig:r2-s2-tests-backend}
\end{figure}


\subsubsection{Analyse de la qualité du code (SonarQube)}

SonarQube analyse le code source backend (\texttt{backend/src}) et frontend
(\texttt{frontend}). La figure~\ref{fig:r2-s2-sonar-conn} présente la configuration de la
connexion entre SonarQube et Azure DevOps. La figure~\ref{fig:r2-s2-sonar-results} illustre
les résultats d'analyse du projet \texttt{ticket-manager}.

\begin{figure}[htbp]
    \centering
    \includegraphics[width=0.85\textwidth]{s2-sonar-connection.png}
    \caption{Connexion SonarQube — Azure DevOps}
    \label{fig:r2-s2-sonar-conn}
\end{figure}

\begin{figure}[htbp]
    \centering
    \includegraphics[width=0.85\textwidth]{s2-sonar-results.png}
    \caption{Résultats d'analyse SonarQube — projet TicketManager}
    \label{fig:r2-s2-sonar-results}
\end{figure}


\subsubsection{Construction des images Docker}

Deux images sont produites : \textbf{ticket-backend} (API NestJS, port 3001) et
\textbf{ticket-frontend} (Next.js, port 3000). La figure~\ref{fig:r2-s2-dockerfile-be}
présente le Dockerfile multi-étapes du backend. La figure~\ref{fig:r2-s2-dockerfile-fe}
détaille celui du frontend. La figure~\ref{fig:r2-s2-docker-build} montre la construction
réussie des images dans le stage Docker du pipeline.

\begin{figure}[htbp]
    \centering
    \includegraphics[width=0.85\textwidth]{s2-dockerfile-backend.png}
    \caption{Dockerfile du backend NestJS}
    \label{fig:r2-s2-dockerfile-be}
\end{figure}

\begin{figure}[htbp]
    \centering
    \includegraphics[width=0.85\textwidth]{s2-dockerfile-frontend.png}
    \caption{Dockerfile du frontend Next.js}
    \label{fig:r2-s2-dockerfile-fe}
\end{figure}

\begin{figure}[htbp]
    \centering
    \includegraphics[width=0.85\textwidth]{s2-docker-build.png}
    \caption{Construction des images dans le stage Docker}
    \label{fig:r2-s2-docker-build}
\end{figure}


\subsubsection{Scan de vulnérabilités (Trivy)}

Trivy analyse les images construites et signale les CVE par niveau de gravité. La
figure~\ref{fig:r2-s2-trivy-be} présente le rapport de scan pour l'image backend. La
figure~\ref{fig:r2-s2-trivy-fe} illustre celui de l'image frontend.

\begin{figure}[htbp]
    \centering
    \includegraphics[width=0.85\textwidth]{s2-trivy-backend.png}
    \caption{Scan Trivy — image backend}
    \label{fig:r2-s2-trivy-be}
\end{figure}

\begin{figure}[htbp]
    \centering
    \includegraphics[width=0.85\textwidth]{s2-trivy-frontend.png}
    \caption{Scan Trivy — image frontend}
    \label{fig:r2-s2-trivy-fe}
\end{figure}


\subsubsection{Publication vers Nexus Repository}

Nexus joue le rôle de \textbf{registre Docker privé} sur vm-gestion. La
figure~\ref{fig:r2-s2-nexus-repo} présente la création du dépôt Docker hosted. La
figure~\ref{fig:r2-s2-nexus-images} montre les images \texttt{ticket-backend} et
\texttt{ticket-frontend} stockées après exécution du pipeline.

\begin{figure}[htbp]
    \centering
    \includegraphics[width=0.85\textwidth]{s2-nexus-repo.png}
    \caption{Création du dépôt Docker hosted dans Nexus}
    \label{fig:r2-s2-nexus-repo}
\end{figure}

\begin{figure}[htbp]
    \centering
    \includegraphics[width=0.85\textwidth]{s2-nexus-images.png}
    \caption{Images \texttt{ticket-backend} et \texttt{ticket-frontend} dans Nexus}
    \label{fig:r2-s2-nexus-images}
\end{figure}


\subsubsection{Terraform automatisé dans le pipeline}

Le stage \textbf{Terraform IaC} exécute \texttt{terraform init}, \texttt{plan} et
\texttt{apply}. La figure~\ref{fig:r2-s2-tf-pipeline} présente l'exécution réussie de ce
stage dans Azure DevOps. La figure~\ref{fig:r2-s2-pipeline-run} illustre une exécution
complète du pipeline CI (stages Build à Push).

\begin{figure}[htbp]
    \centering
    \includegraphics[width=0.85\textwidth]{s2-terraform-pipeline.png}
    \caption{Exécution réussie du stage Terraform dans Azure DevOps}
    \label{fig:r2-s2-tf-pipeline}
\end{figure}

\begin{figure}[htbp]
    \centering
    \includegraphics[width=0.85\textwidth]{s2-pipeline-run-success.png}
    \caption{Exécution complète du pipeline CI (stages Build à Push)}
    \label{fig:r2-s2-pipeline-run}
\end{figure}


% ---------------------------------------------------------------------------
\subsection{Bilan du Sprint 2}
% ---------------------------------------------------------------------------

Le Sprint~2 met en place une chaîne d'intégration continue couvrant la compilation, la
qualité, la conteneurisation, la sécurité des images et la publication vers un registre
privé. Le déploiement effectif des pods applicatifs sur le cluster (stage
\textbf{DeployK8s}) et la supervision (stage \textbf{DeployMonitoring}) sont détaillés au
Sprint~3.


% --- Sprint 3 : Monitoring et CD ----------------------------------------------
% =============================================================================
% Release 2 — Sprint 3 : Monitoring et supervision de la plateforme
% =============================================================================

\section[Sprint 3 : Monitoring et supervision de la plateforme]%
{Sprint 3 : Monitoring et supervision de la plateforme}

Ce sprint finalise la Release~2 par le \textbf{déploiement continu} (CD) de l'application
sur Kubernetes et la mise en place de la \textbf{supervision} (Prometheus, Grafana) via
Ansible sur vm-gestion.


% ---------------------------------------------------------------------------
\subsection{Objectif du Sprint 3}
% ---------------------------------------------------------------------------

Le tableau~\ref{tab:r2-s3-objectifs} détaille les objectifs du Sprint~3 par axe.

\begin{table}[H]
\centering
\caption{Objectifs du Sprint~3 par axe}
\label{tab:r2-s3-objectifs}
\footnotesize
\renewcommand{\arraystretch}{1.1}
\setlength{\tabcolsep}{6pt}
\begin{tabular}{|p{2.8cm}|p{9.2cm}|}
\hline
\textbf{Axe} & \textbf{Objectif} \\
\hline
Déploiement continu &
  Déployer automatiquement backend, frontend et MongoDB in-cluster via le stage
  \texttt{DeployK8s} (branche \texttt{main}, approbation manuelle) \\
\hline
Manifests Kubernetes &
  Appliquer les fichiers YAML du répertoire \texttt{k8s/} (Deployments, Services,
  StatefulSet MongoDB, Secrets) \\
\hline
Registre privé &
  Configurer le pull d'images depuis Nexus (\texttt{10.1.1.4:8082}) sur les workers \\
\hline
Supervision &
  Déployer Prometheus et Grafana sur vm-gestion (\texttt{ansible/playbooks/monitoring.yml}) \\
\hline
Observabilité &
  Collecter les métriques des nœuds K8s, du backend NestJS et de MongoDB \\
\hline
Validation &
  Smoke tests post-déploiement (API, frontend, cibles Prometheus UP) \\
\hline
\end{tabular}
\end{table}

\paragraph{Résultat attendu.}
\begin{itemize}
  \item namespace \texttt{ticketmanager} avec pods \texttt{Running} ;
  \item application accessible via NodePort (frontend \texttt{30080}, backend \texttt{30081}) ;
  \item Prometheus accessible sur le port \texttt{9090}, Grafana sur \texttt{3030} ;
  \item dashboards opérationnels pour le cluster et l'application.
\end{itemize}

\paragraph{Choix architectural — MongoDB in-cluster.}
Contrairement aux projets précédents (chatbot) où MongoDB était exclusivement sur une VM
dédiée, TicketManager héberge la base de données de production dans un
\textbf{StatefulSet Kubernetes} avec volume persistant. La vm-db reste disponible pour
migration ou sauvegarde legacy (\texttt{ansible/playbooks/mongodb.yml}).


% ---------------------------------------------------------------------------
\subsection{Backlog du Sprint 3}
% ---------------------------------------------------------------------------

Le tableau~\ref{tab:r2-s3-backlog-deploy} couvre l'epic déploiement continu. Le
tableau~\ref{tab:r2-s3-backlog-monitoring} couvre l'epic supervision.

\subsubsection{Epic — Déploiement continu (Kubernetes)}

\begin{table}[H]
\centering
\caption{Backlog Sprint~3 — Epic déploiement continu Kubernetes}
\label{tab:r2-s3-backlog-deploy}
\footnotesize
\renewcommand{\arraystretch}{0.85}
\setlength{\tabcolsep}{4pt}
\begin{tabular}{|c|p{7.8cm}|c|c|}
\hline
\textbf{ID} & \textbf{User story} & \textbf{Priorité} & \textbf{Est. (j)} \\
\hline
R2-S3-D01 &
  En tant que développeur, je déploie le backend NestJS en pod Kubernetes &
  H & 2 \\
\hline
R2-S3-D02 &
  En tant que développeur, je déploie le frontend Next.js en pod Kubernetes &
  H & 2 \\
\hline
R2-S3-D03 &
  En tant que développeur, je déploie MongoDB in-cluster via un StatefulSet &
  H & 2 \\
\hline
R2-S3-D04 &
  En tant que développeur, je configure les secrets Kubernetes (JWT, MongoDB, Cloudinary) &
  H & 1 \\
\hline
R2-S3-D05 &
  En tant que développeur, j'automatise le déploiement via le stage CD du pipeline &
  H & 2 \\
\hline
R2-S3-D06 &
  En tant que développeur, je valide le déploiement par des smoke tests automatisés &
  H & 1 \\
\hline
\end{tabular}
\end{table}

\subsubsection{Epic — Monitoring et supervision}

\begin{table}[H]
\centering
\caption{Backlog Sprint~3 — Epic monitoring et supervision}
\label{tab:r2-s3-backlog-monitoring}
\footnotesize
\renewcommand{\arraystretch}{0.85}
\setlength{\tabcolsep}{4pt}
\begin{tabular}{|c|p{7.8cm}|c|c|}
\hline
\textbf{ID} & \textbf{User story} & \textbf{Priorité} & \textbf{Est. (j)} \\
\hline
R2-S3-M01 &
  En tant que développeur, je déploie Prometheus sur vm-gestion via Ansible &
  H & 2 \\
\hline
R2-S3-M02 &
  En tant que développeur, je déploie Grafana et je connecte Prometheus comme source de données &
  H & 2 \\
\hline
R2-S3-M03 &
  En tant que développeur, j'intègre le déploiement monitoring au pipeline Azure DevOps &
  M & 1 \\
\hline
R2-S3-M04 &
  En tant que développeur, je vérifie que les cibles Prometheus sont en état \texttt{UP} &
  H & 1 \\
\hline
\end{tabular}
\end{table}


% ---------------------------------------------------------------------------
\subsection{Analyse — Conception du déploiement Kubernetes}
% ---------------------------------------------------------------------------

\subsubsection{Architecture de déploiement}

Les composants applicatifs sont encapsulés dans des pods du namespace
\texttt{ticketmanager} :
\begin{itemize}
  \item \textbf{ticket-backend} — image Nexus, port 3001, secrets JWT et Cloudinary ;
  \item \textbf{ticket-frontend} — image Nexus, port 3000, URL API figée au build Docker ;
  \item \textbf{mongodb} — StatefulSet avec volume persistant pour les données tickets.
\end{itemize}

Les requêtes externes atteignent les services via \textbf{NodePort} sur les workers
(\texttt{30080} frontend, \texttt{30081} backend). La figure~\ref{fig:r2-s3-k8s-arch}
détaille l'architecture de déploiement Kubernetes (ressources du namespace
\texttt{ticketmanager}), en complément de la vue d'ensemble physique (figure~2.5).

\begin{figure}[htbp]
    \centering
    \includegraphics[width=0.92\textwidth]{architecture-deploiement-k8s.png}
    \caption{Architecture de déploiement Kubernetes — namespace \texttt{ticketmanager}}
    \label{fig:r2-s3-k8s-arch}
\end{figure}

\subsubsection{Chaîne de déploiement continu}

Le stage \texttt{DeployK8s} du pipeline :
\begin{enumerate}
  \item attend une \textbf{approbation manuelle} (\texttt{ManualValidation}) ;
  \item configure le registre insecure sur les workers ;
  \item exécute \texttt{k8s/deploy-pipeline.sh} (secrets, manifests, rollout, seed, smoke test).
\end{enumerate}


% ---------------------------------------------------------------------------
\subsection{Réalisation — Déploiement Kubernetes}
% ---------------------------------------------------------------------------

\subsubsection{Manifests et secrets}

La figure~\ref{fig:r2-s3-backend-yaml} présente le manifest de déploiement du backend
(\texttt{backend-deployment.yaml}). La figure~\ref{fig:r2-s3-frontend-yaml} illustre celui
du frontend (\texttt{frontend-deployment.yaml}).

\begin{figure}[htbp]
    \centering
    \includegraphics[width=0.85\textwidth]{s3-backend-deployment.png}
    \caption{Manifest de déploiement du backend (\texttt{backend-deployment.yaml})}
    \label{fig:r2-s3-backend-yaml}
\end{figure}

\begin{figure}[htbp]
    \centering
    \includegraphics[width=0.85\textwidth]{s3-frontend-deployment.png}
    \caption{Manifest de déploiement du frontend (\texttt{frontend-deployment.yaml})}
    \label{fig:r2-s3-frontend-yaml}
\end{figure}

\subsubsection{Exécution du déploiement}

La figure~\ref{fig:r2-s3-k8s-pipeline} présente le stage \textbf{DeployK8s} exécuté dans
Azure DevOps, incluant l'approbation manuelle et l'application des manifests.

\begin{figure}[htbp]
    \centering
    \includegraphics[width=0.85\textwidth]{s3-k8s-pipeline-stage.png}
    \caption{Stage DeployK8s dans Azure DevOps}
    \label{fig:r2-s3-k8s-pipeline}
\end{figure}

\subsubsection{Vérification des pods}

La figure~\ref{fig:r2-s3-pods} présente la sortie de \texttt{kubectl get pods -n ticketmanager}
: l'ensemble des pods applicatifs est en état \texttt{Running}.

\begin{figure}[htbp]
    \centering
    \includegraphics[width=0.85\textwidth]{s3-kubectl-pods.png}
    \caption{État des pods (\texttt{kubectl get pods -n ticketmanager})}
    \label{fig:r2-s3-pods}
\end{figure}

\subsubsection{Application accessible}

La figure~\ref{fig:r2-s3-app} illustre l'interface TicketManager accessible depuis un
navigateur après déploiement réussi sur le cluster.

\begin{figure}[htbp]
    \centering
    \includegraphics[width=0.85\textwidth]{s3-app-browser.png}
    \caption{Interface TicketManager accessible après déploiement}
    \label{fig:r2-s3-app}
\end{figure}


% ---------------------------------------------------------------------------
\subsection{Réalisation — Monitoring et supervision}
% ---------------------------------------------------------------------------

\subsubsection{Déploiement Ansible (Prometheus + Grafana)}

Le playbook \texttt{ansible/playbooks/monitoring.yml} déploie Prometheus (port 9090) et
Grafana (port 3030) sur vm-gestion. La figure~\ref{fig:r2-s3-ansible} présente l'exécution
réussie de ce playbook.

\begin{figure}[htbp]
    \centering
    \includegraphics[width=0.85\textwidth]{s3-ansible-monitoring.png}
    \caption{Exécution du playbook \texttt{monitoring.yml}}
    \label{fig:r2-s3-ansible}
\end{figure}

\subsubsection{Configuration Prometheus}

Le fichier \texttt{ansible/templates/prometheus.yml.j2} définit les cibles surveillées.
La figure~\ref{fig:r2-s3-prom-targets} présente l'interface Prometheus (\texttt{/targets})
avec l'ensemble des cibles en état \texttt{UP}.

\begin{figure}[htbp]
    \centering
    \includegraphics[width=0.85\textwidth]{s3-prometheus-targets.png}
    \caption{Cibles actives dans l'interface Prometheus (\texttt{/targets})}
    \label{fig:r2-s3-prom-targets}
\end{figure}

\subsubsection{Dashboards Grafana}

La figure~\ref{fig:r2-s3-grafana-ds} présente la configuration de Prometheus comme source
de données dans Grafana. La figure~\ref{fig:r2-s3-grafana-dash} illustre un dashboard de
supervision regroupant les métriques du cluster et de l'application backend.

\begin{figure}[htbp]
    \centering
    \includegraphics[width=0.85\textwidth]{s3-grafana-datasource.png}
    \caption{Configuration de Prometheus comme source Grafana}
    \label{fig:r2-s3-grafana-ds}
\end{figure}

\begin{figure}[htbp]
    \centering
    \includegraphics[width=0.85\textwidth]{s3-grafana-dashboard.png}
    \caption{Dashboard de supervision — cluster et application}
    \label{fig:r2-s3-grafana-dash}
\end{figure}

\subsubsection{Intégration pipeline}

La figure~\ref{fig:r2-s3-mon-pipeline} présente le stage \textbf{DeployMonitoring} du
pipeline Azure DevOps, exécutant automatiquement le playbook Ansible après un déploiement
K8s réussi.

\begin{figure}[htbp]
    \centering
    \includegraphics[width=0.85\textwidth]{s3-monitoring-pipeline.png}
    \caption{Stage DeployMonitoring dans Azure DevOps}
    \label{fig:r2-s3-mon-pipeline}
\end{figure}


% ---------------------------------------------------------------------------
\subsection{Tests et validation}
% ---------------------------------------------------------------------------

\subsubsection{Smoke tests post-déploiement}

Le script \texttt{k8s/post-deploy.sh} vérifie la disponibilité de l'API, du frontend et
des outils de supervision. Le tableau~\ref{tab:r2-s3-validation} résume les critères de
validation retenus pour le Sprint~3.

\begin{table}[H]
\centering
\caption{Critères de validation du Sprint~3}
\label{tab:r2-s3-validation}
\footnotesize
\renewcommand{\arraystretch}{1.1}
\setlength{\tabcolsep}{6pt}
\begin{tabular}{|l|l|l|}
\hline
\textbf{Composant} & \textbf{URL / commande} & \textbf{Critère de succès} \\
\hline
Backend & \texttt{curl http://<worker>:30081/} & HTTP 200 \\
\hline
Frontend & \texttt{http://<worker>:30080} & Page de connexion affichée \\
\hline
Pods & \texttt{kubectl get pods -n ticketmanager} & Tous \texttt{Running} \\
\hline
Prometheus & \texttt{curl http://10.1.1.4:9090/-/ready} & \texttt{Ready} \\
\hline
Grafana & \texttt{http://10.1.1.4:3030/login} & HTTP 200 \\
\hline
\end{tabular}
\end{table}


% ---------------------------------------------------------------------------
\subsection{Bilan du Sprint 3}
% ---------------------------------------------------------------------------

Le Sprint~3 clôt la Release~2 en rendant TicketManager déployable et observable en
environnement cloud. L'application tourne dans Kubernetes ; la supervision Prometheus/Grafana
permet de détecter les anomalies de charge, de disponibilité et de performance avant
intervention manuelle.


% --- Bilan Release 2 ----------------------------------------------------------
\subsection{Bilan de la Release~2}

À l'issue des trois sprints, la plateforme \textit{TicketManager} dispose d'une
infrastructure documentée et reproductible sur Azure Stack Hub, d'une chaîne CI/CD
automatisée (compilation, qualité, conteneurisation, publication Nexus, provisionnement
IaC) et d'un déploiement applicatif supervisé (pods Kubernetes, Prometheus, Grafana).

Les indicateurs de réussite retenus pour cette release sont :
\begin{itemize}
  \item exécution réussie du pipeline Azure DevOps sur la branche \texttt{main} ;
  \item pods du namespace \texttt{ticketmanager} en état \texttt{Running} ;
  \item accessibilité de l'interface web et de l'API depuis le réseau autorisé ;
  \item cibles Prometheus en état \texttt{UP} et dashboards Grafana opérationnels.
\end{itemize}

Les évolutions ultérieures (hors Release~2) peuvent porter sur le renforcement des tests
automatisés dans le pipeline, la réactivation des stages Snyk/Trivy, l'exposition via
Ingress Controller et la haute disponibilité multi-réplicas.
